%% =====================================================================
%% PAER-LLM -- IEEE/ACM Transactions on Audio, Speech, and Language
%% Processing (T-ASLP) submission.
%%
%% Compile:  pdflatex paer_llm_ieee -> bibtex paer_llm_ieee -> pdflatex x2
%%
%% Requires: references.bib
%%           figures/fig1_framework_ieee.pdf         (supplied)
%%           figures/fig4_protocol_vs_architecture.pdf (supplied)
%%           figures/fig8_confusion_matrices.pdf       (supplied)
%%           figures/fig9_minimum_detectable_effect.pdf(supplied)
%%           figures/fig2_decomposition_ladder.pdf   <- to produce
%%           figures/fig5_stream_complementarity.pdf <- to produce
%%           figures/fig7_reliability.pdf            <- to produce
%%
%% Every numerical value below was recomputed from the project's own
%% result files:
%%   Evaluation/cv_results_{concat,gated,cross_attention}_speaker.json
%%   Evaluation/test_metrics.json
%%   Evaluation/test_metrics_cross_attention_random_dup.json
%%   Evaluation/classification_report*.txt, confusion_matrix*.csv
%%   Models/train_config*.json, Models/best_model.config.json
%%
%% DRAFT MARKERS: see the \ifdraft block below.
%% LAYOUT: review (one column) vs camera-ready (two column) -- one line.
%% =====================================================================

% ---------------------------------------------------------------------
% LAYOUT. Exactly one of these must be active.
%
%   review     one column, 11pt, generous leading. This is IEEE's own
%              review-copy format: broader measure, more air, easier to
%              annotate. Roughly twice the page count of camera-ready,
%              which is expected and not a problem at submission.
%   cameraready  the two-column T-ASLP page.
%
% Everything downstream adapts automatically -- figure widths, wide
% floats and table shrink-to-fit all key off \ifCLASSOPTIONonecolumn.
% ---------------------------------------------------------------------
\documentclass[journal,onecolumn,11pt,draftclsnofoot]{IEEEtran}   % review
%\documentclass[journal]{IEEEtran}                                % camera-ready

\usepackage{amsmath,amssymb,amsthm}
\usepackage{bm}
\usepackage{booktabs}
\usepackage{graphicx}
\usepackage{multirow}
\usepackage{array}
\newcolumntype{L}[1]{>{\raggedright\arraybackslash}p{#1}}
\usepackage{algorithm}
\usepackage{algpseudocode}
\usepackage{xcolor}
\usepackage{url}
\usepackage[hidelinks]{hyperref}
\usepackage{cite}
\usepackage{adjustbox}
\usepackage{setspace}

% ---------------------------------------------------------------------
% Layout-dependent measures.
%
% In one-column mode \columnwidth equals \textwidth, so a figure drawn
% for a 3.4in column would be blown up to 6.5in and its type would swell
% with it. \figw keeps every figure at the size it was designed for in
% either mode.
% ---------------------------------------------------------------------
\newlength{\figw}
\newlength{\wfigw}
\ifCLASSOPTIONonecolumn
  \setstretch{1.28}          % airy, but short of full double spacing
  \AtBeginDocument{%
    \setlength{\figw}{0.62\columnwidth}%
    \setlength{\wfigw}{0.86\textwidth}}
\else
  \AtBeginDocument{%
    \setlength{\figw}{\columnwidth}%
    \setlength{\wfigw}{0.95\textwidth}}
\fi

% ---------------------------------------------------------------------
% Draft markers.
%
%   \drafttrue   working copy: every gap shouts in red, with the command
%                that fills it.
%   \draftfalse  clean copy: prose markers disappear and unfilled result
%                cells become a discreet grey "pending".
%
% \draftfalse deliberately does NOT print "---" for an unfilled cell. A
% dash reads as a legitimate "not applicable", so the one failure mode
% this paper exists to criticise -- a missing number that looks like a
% reported one -- would be reintroduced by its own template. Before
% submitting, check that nothing is left:
%
%     grep -n "TORUN\|TOADD\|\\RUN\b" paer_llm_ieee.tex
% ---------------------------------------------------------------------
\newif\ifdraft
\draftfalse                   % \drafttrue while filling in results
\ifdraft
  \newcommand{\TORUN}[1]{\textcolor{red}{\textbf{[run: #1]}}}
  \newcommand{\TOADD}[1]{\textcolor{red}{\textbf{[add: #1]}}}
  \newcommand{\RUN}{\textcolor{red}{\textbf{--?--}}}
\else
  % \TORUN is a note to the author sitting at the end of a sentence that
  % reads correctly without it, so it simply goes away.
  \newcommand{\TORUN}[1]{}
  % \TOADD marks a hole *inside* a sentence -- the affiliation, the
  % repository URL. Dropping it silently leaves "G. Sharma is with
  % (e-mail: ...)", so in clean mode it stays, quietly.
  \newcommand{\TOADD}[1]{\textcolor{gray}{[#1]}}
  \newcommand{\RUN}{\textcolor{gray}{\footnotesize pending}}
\fi

% ---------------------------------------------------------------------
% \paperfig{path-without-extension}{width}{reserved height}
%
% Draws the graphic if the file is on disk, and otherwise reserves an
% area of exactly the height the finished figure will occupy. Running
% Experiments/make_paper_figures.py is therefore the whole of "insert
% the figures" -- there is no LaTeX edit to remember afterwards, and no
% window in which the document fails to compile.
% ---------------------------------------------------------------------
\newcommand{\reservedfig}[2]{%
  \fboxsep=0pt\fboxrule=0.6pt%
  \textcolor{gray}{\fbox{\parbox[c][#2][c]{#1}{%
    \hbadness=10000\relax\centering\tiny\itshape\color{gray}
    reserved space\\ figure to be inserted\\[1pt]
    (file named in the caption)}}}}

% #3 is the finished figure's aspect ratio (height/width), so the space
% reserved is the space the graphic will actually occupy in whichever
% layout is active -- the page never reflows when the figures arrive.
\newlength{\pfw}
\newcommand{\paperfig}[3]{%
  \IfFileExists{#1.pdf}%
    {\includegraphics[width=#2]{#1.pdf}}%
    {\setlength{\pfw}{#2}\reservedfig{\pfw}{#3\pfw}}}

\newtheorem{proposition}{Proposition}
\newtheorem{corollary}{Corollary}
\newtheorem{definition}{Definition}

\begin{document}

\title{A Unified Psychoacoustic--Semantic Representation for Speech
Emotion Recognition, and the Evaluation Protocol Required to Assess It}

\author{Gaurav~Sharma%
\thanks{G. Sharma is with \TOADD{Department, Institution, City, Postcode,
Country} (e-mail: gaurav.vit08@gmail.com).}%
\thanks{Manuscript submitted \TOADD{date}.}}

\markboth{IEEE/ACM Transactions on Audio, Speech, and Language Processing,
Vol.~XX, No.~X, XXXX}{Sharma: A Unified Psychoacoustic--Semantic
Representation for Speech Emotion Recognition}

\IEEEtitleabstractindextext{%
\begin{abstract}
Speech emotion recognition increasingly relies on self-supervised
representations, yet whether perceptually motivated psychoacoustic descriptors
remain complementary to them is unresolved, and published eight-class
accuracies on the RAVDESS corpus range from 39.6\% to 99.0\%. We propose
PAER-LLM, a five-layer framework that unifies 34 psychoacoustic descriptors
with frozen HuBERT embeddings through a learned per-dimension gate and
conditions an instruction-tuned language model on the resulting posterior. We
formalise the unified-representation problem, the protocol-dependence of its
assessment, and the resolution at which a representational claim is
admissible. Evaluating one implementation under six protocols that differ only
in partitioning and model selection, with architecture, features, optimiser and
schedule held fixed, accuracy falls from 94.79\% under a contaminated random
split to 64.28\%~$\pm$~5.30 under speaker-disjoint five-fold
cross-validation---a spread of 30.5 points. Under the most defensible protocol
the largest difference between three fusion operators is 2.03 points, below the
design's 3.47-point minimum detectable effect. Contamination also conceals a
21.6-point gender disparity, shrinking it to 2.5 points. We prove that
single-token cross-modal attention is analytically degenerate, its query stream
receiving exactly zero gradient, and show that a corrected multi-token form is
outperformed by a single gate at 3.1$\times$ the parameters.
\end{abstract}

\begin{IEEEkeywords}
Speech emotion recognition, psychoacoustic features, representation learning,
feature fusion, evaluation protocol, data leakage, speaker-independent
evaluation, large language models.
\end{IEEEkeywords}}

\maketitle
\IEEEdisplaynontitleabstractindextext
\IEEEpeerreviewmaketitle

% =====================================================================
\section{Introduction}
\label{sec:intro}
% =====================================================================
\IEEEPARstart{S}{peech} emotion recognition (SER) underpins affective
human--computer interaction, and is supported by large self-supervised speech
encoders and an extensive benchmark literature \cite{wani2021comprehensive,
latif2023survey,khare2024emotion}. This paper is concerned with two problems
that the field treats separately and that, we argue, cannot be separated.

\subsection{Problem Statement}
\label{sec:problem}

Emotion in speech is carried by perceptually salient acoustic properties.
Listeners recover a talker's affective state from fundamental frequency and its
movement, from intensity, and from how energy is distributed across the
auditory critical bands---quantities that hearing science has described for six
decades and that standardised parameter sets such as GeMAPS/eGeMAPS formalise
\cite{eyben2016gemaps,zwicker1961subdivision}. Modern SER, however, is built
almost entirely on self-supervised transformer encoders trained on raw
waveforms with no perceptual objective \cite{baevski2020wav2vec2,hsu2021hubert,
chen2022wavlm}. These two accounts of what matters in an emotional utterance
have largely stopped talking to each other. The first problem this paper
addresses is therefore representational: \emph{do perceptually motivated
psychoacoustic descriptors still carry emotional information that a large
self-supervised encoder does not already provide, and if so, how should the two
be combined into a single representation?} The question is not settled. Dixit
\emph{et al.} \cite{dixit2024explaining} regress self-supervised embeddings onto
eGeMAPS categories and find that deep representations recover only part of the
handcrafted perceptual space, which implies complementarity; yet systems that
add handcrafted descriptors to self-supervised embeddings report gains of one
to three accuracy points, which is small enough that it may not be a gain at
all.

That last clause is the second problem, and it turns out to govern the first.
Published eight-class accuracies on the Ryerson Audio-Visual Database of
Emotional Speech and Song (RAVDESS) \cite{livingstone2018ravdess}---the corpus
on which such comparisons are most often made---span 39.6\%
\cite{goel2024camulenet} to 99.0\% \cite{jafarzadeh2024ssl} within the last two
years. They cannot both characterise the same task. The discrepancy is not
principally architectural: it follows from how the corpus is partitioned.
RAVDESS comprises 1440 speech utterances produced by 24 professional actors,
each reading the same two sentences across all eight emotion categories at two
intensities with two repetitions. A random utterance-level partition therefore
guarantees that for every test item there exists a training item from the same
speaker, with the same lexical content, at the same emotional intensity,
differing only in take. A classifier can succeed by recognising the talker
rather than the emotion. When the confounding effect of the evaluation design
exceeds the effect being measured, no representational question can be answered
on that evidence, however carefully the representation is built.

The two problems compound. A researcher proposing a unified psychoacoustic and
self-supervised representation, as we do, has no way to establish that the
psychoacoustic stream earns its place if the measurement apparatus cannot
resolve differences of the size that stream is likely to produce. Conversely, a
field that reports architectural improvements of one to three points under
undisclosed partitioning schemes accumulates a literature in which
representational claims cannot be checked. This paper therefore treats the two
as one problem. We formalise the unified representation and the protocol that
assesses it in a single framework (Section~\ref{sec:formulation}), construct the
representation (Section~\ref{sec:framework}), state precisely how much evidence
would be required to accept it, and then measure how much evidence the field's
standard protocols actually supply. The answer is: considerably less than is
required, and the shortfall is large enough to swallow the entire architectural
literature on this corpus.

\subsection{Research Objective and Scope}
\label{sec:objective}

The objective of this work is to propose and characterise a unified
representation framework that integrates psychoacoustic emotional features with
self-supervised speech representations, and to condition a large language model
on the emotional evidence that framework produces. We name the framework
PAER-LLM (Psychoacoustically-Aware Emotion Recognition with a Language Model)
and specify it as five layers in Section~\ref{sec:framework}.

Three boundaries on that objective should be stated at the outset.

First, the framework is \emph{unimodal}. Both streams are computed from one
audio signal; they are two feature views of the same modality, not two
modalities. We therefore use ``unified representation'' and never
``multimodal'', and distinguish our setting from audio-visual work explicitly.

Second, Layers~1--3 are evaluated here. Layer~4, which conditions an
instruction-tuned language model on the classifier's posterior, is specified
and implemented but \emph{not} evaluated: we report no human study, no
automatic metric and no ablation for it. Layer~5 is a demonstration interface.
We mark this division in Fig.~\ref{fig:framework} rather than leaving it to be
inferred, and Section~\ref{sec:llm-discussion} sets out the calibration
experiment that would make Layer~4 load-bearing.

Third, the empirical scope is a single corpus. All numerical results are on
RAVDESS. The analytical results of Section~\ref{sec:properties} are
corpus-independent; the magnitudes are not.

We stress one point of interpretation. The configuration of PAER-LLM that we
adopt is statistically indistinguishable from plain concatenation of the two
encoded streams, and its originally proposed fusion mechanism does not work at
all, in a sense we prove in Proposition~\ref{prop:degen}. We report both
against our own interest, because the alternative---selecting the configuration
that flatters the proposed method---is precisely the researcher degree of
freedom that Nie{\ss}l \emph{et al.} \cite{niessl2022overoptimism} identify as a
source of systematic over-optimism.

\subsection{Contributions}
\label{sec:contributions}

\begin{enumerate}
  \item \textbf{A formal problem formulation}
    (Section~\ref{sec:formulation}) that states the unified-representation
    objective, the information-theoretic condition under which the
    psychoacoustic view can contribute, the evaluation protocol as an operator
    on the corpus, and an \emph{admissibility} criterion that a representational
    claim must satisfy to be supported by a given design.
  \item \textbf{A unified psychoacoustic--semantic representation framework.}
    PAER-LLM encodes 34 perceptually motivated descriptors and a frozen
    mean-pooled HuBERT-base embedding into a common 128-dimensional space and
    combines them with a learned per-dimension gate, whose posterior conditions
    an instruction-tuned language model (Section~\ref{sec:framework},
    Algorithms~\ref{alg:unified}--\ref{alg:prompt}). All three fusion operators
    compared emit exactly 128 dimensions, so downstream capacity is fixed and
    the comparison isolates the operator.
  \item \textbf{Three analytical properties of pooled-vector fusion}
    (Section~\ref{sec:properties}): single-token cross-modal attention is
    exactly degenerate and gradient-inert; gated fusion is confined to the
    per-coordinate convex hull of its inputs while concatenation is not, so
    neither operator contains the other; and the parameter and multiply--
    accumulate cost of token cross-attention exceeds both by 3.1$\times$ and
    5.3$\times$ respectively.
  \item \textbf{A controlled decomposition of the evaluation-protocol effect}
    (Section~\ref{sec:protocol}, Algorithm~\ref{alg:decompose}), isolating
    duplicate records, speaker-dependent partitioning, test-set model selection
    and single-split reporting within one fixed system. To our knowledge this is
    the first study to attribute a RAVDESS performance gap to its individual
    protocol components rather than to report two endpoints.
  \item \textbf{A protocol-to-architecture effect ratio and a stated resolution
    limit.} Under the most defensible protocol the largest difference between
    three fusion operators is 2.03 points against a 30.5-point protocol effect;
    the design's minimum detectable effect is 3.47 points, so the entire
    architectural range is inadmissible under our own criterion
    (Sections~\ref{sec:fusion} and~\ref{sec:mde}).
  \item \textbf{Evidence that contamination conceals fairness failures.} Under
    the contaminated protocol the female--male accuracy disparity is 2.5 points;
    under speaker-disjoint evaluation of the same system it is 21.6 points, and
    the dominant per-class failure (\emph{happy}, $F_1$ 0.511) is invisible
    ($F_1$ 0.915). Protocol choice therefore distorts not only reported accuracy
    but the fairness and error audits built on it (Section~\ref{sec:speaker}).
  \item \textbf{A reproducible evaluation harness} with published fold
    assignments, the contaminated-protocol reproduction, a corpus-integrity
    check, the complementarity and calibration analysis, and per-fold prediction
    files.
\end{enumerate}

\subsection{Organisation}
Section~\ref{sec:related} reviews the relevant literature including recent
systems. Section~\ref{sec:formulation} formalises the problem.
Section~\ref{sec:framework} specifies the framework.
Section~\ref{sec:methodology} sets out the research methodology and the
benchmarking design. Section~\ref{sec:setup} gives the experimental
configuration. Section~\ref{sec:results} reports results, including the
performance matrix and complexity analysis. Section~\ref{sec:discussion}
discusses, Section~\ref{sec:limitations} states limitations, and
Section~\ref{sec:conclusion} concludes.

% =====================================================================
\section{Related Work}
\label{sec:related}
% =====================================================================

\subsection{Self-Supervised Representations for Emotion Recognition}
Self-supervised encoders---wav2vec~2.0 \cite{baevski2020wav2vec2}, HuBERT
\cite{hsu2021hubert} and WavLM \cite{chen2022wavlm}---underpin most competitive
SER systems, and SUPERB \cite{yang2021superb} established frozen-encoder probing
as a standard evaluation. Pepino \emph{et al.} \cite{pepino2021wav2vec} showed
that a learned weighted combination of intermediate layers substantially
outperforms the final layer for emotion, consistent with layer-wise probing
analyses showing an acoustic--phonetic--acoustic trajectory across depth
\cite{pasad2021layerwise,pasad2023comparative}. Chiu \emph{et al.}
\cite{chiu2025probing} report that speaker-state and prosodic attributes are
most robustly represented in middle layers while final layers suppress speaker
identity. Wagner \emph{et al.} \cite{wagner2023dawn} provide the most systematic
study of transformer fine-tuning for dimensional SER, including robustness and
fairness analyses. Emotion-specialised pretraining extends this line
\cite{chen2024vesper,ma2024emotion2vec}.

We use frozen mean-pooled HuBERT-base representations. This is deliberately
conservative: it isolates the contribution of the psychoacoustic stream and the
fusion operator from that of encoder adaptation, which is known to be large
\cite{chen2023exploring}.

\subsection{Psychoacoustic and Prosodic Representations}
That emotion is carried by perceptually salient acoustic correlates remains
supported by recent systematic review \cite{schewski2025acoustic,
larrouymaestri2025prosody}. Standardised parameter sets such as GeMAPS/eGeMAPS
\cite{eyben2016gemaps}, extracted via openSMILE \cite{eyben2010opensmile},
formalise this premise. The critical-band concept and its analytical
approximation are due to Zwicker \cite{zwicker1961subdivision} and Zwicker and
Terhardt \cite{zwicker1980analytical} respectively. Whether such descriptors
retain value alongside self-supervised embeddings is the representational
question this framework is built to test \cite{dixit2024explaining}.
Perceptually motivated spectral front-ends have shown advantages over mel
filterbanks, particularly under noise \cite{peng2023cochleagram,
peng2023dimensional}.

\subsection{Fusion of Heterogeneous Feature Streams}
Cross-attention is the default mechanism for fusing heterogeneous
representations in affective computing \cite{vaswani2017attention,
praveen2021crossattention,zou2022coattention}, including systems fusing
handcrafted acoustic descriptors with self-supervised embeddings
\cite{zhao2025humpcat}. Gated variants are common \cite{li2025wavfusion}, and
Nasersharif and Azad \cite{nasersharif2024attention} compare attention-based
fusion against simpler classifier combination. Much of this literature is
genuinely multimodal---audio with video or text---whereas our two streams are
views of one signal; the distinction matters because cross-modal attention is
motivated by alignment between sequences that our pooled setting does not
provide (Section~\ref{sec:why-fusion}).

A parallel literature questions how much attention contributes. Dong \emph{et
al.} \cite{dong2021attentionnotall} prove that pure attention loses rank doubly
exponentially with depth; Zhai \emph{et al.} \cite{zhai2023entropycollapse}
characterise pathological attention-entropy regimes; and Bibal \emph{et al.}
\cite{bibal2022attention} survey whether attention weights constitute evidence
of learned interaction. Proposition~\ref{prop:degen} is related but distinct:
rank collapse is asymptotic in depth, whereas our degeneracy is \emph{exact} and
arises from sequence length one, independently of depth, head count or scaling.

\subsection{Language Models for Emotion-Sensitive Interaction}
Emotionally aware dialogue has been studied through emotional support
conversation \cite{liu2021esconv}, empathetic response generation
\cite{qian2023empathetic}, and benchmarks of emotional competence
\cite{sabour2024emobench}. Architecturally, systems divide into end-to-end
audio-conditioned models \cite{tang2023salmonn,chu2024qwen2audio,
kim2024paralinguistics} and cascaded designs. Recent work shows a frozen LLM can
be made to perceive paralinguistic style through encoder alignment
\cite{kang2025frozen}, and that structured, interpretable speech descriptors can
condition an LLM for explainable emotion recognition \cite{chen2025descriptors}.
We adopt a cascaded design conditioned on the classifier's full posterior rather
than its top-1 label, informed by work on eliciting calibrated confidence
\cite{tian2023calibration} and on principled abstention \cite{wen2025abstention}.

\subsection{Evaluation Protocol, Leakage and Reproducibility}
\label{sec:rw-leakage}
Kapoor and Narayanan \cite{kapoor2023leakage} document leakage-driven
irreproducibility across seventeen fields. Elangovan \emph{et al.}
\cite{elangovan2021memorization} quantify train--test overlap in NLP benchmarks
and show that near-duplicate examples permit high scores through memorisation.
Rosenblatt \emph{et al.} \cite{rosenblatt2024leakage} find that repeated
subjects across partitions inflate performance most severely, with small
datasets most vulnerable---a description that fits RAVDESS closely. Gorman and
Bedrick \cite{gorman2019splits} show that conclusions drawn from standard splits
frequently fail to replicate under resampling, and Bouthillier \emph{et al.}
\cite{bouthillier2021variance} demonstrate that ignoring sources of benchmark
variance inflates apparent significance \cite{niessl2022overoptimism,
lones2024pitfalls}.

The standard we advocate is not new. The INTERSPEECH 2009 Emotion Challenge
\cite{schuller2009challenge} established speaker-independent partitioning as the
norm nearly two decades ago, and the Interspeech 2025 challenge on naturalistic
SER \cite{naini2025challenge} retains it. EmoBox \cite{ma2024emobox} now
distributes standardised speaker-independent splits across corpora, including
RAVDESS. What has happened since 2009 is erosion rather than absence.

\subsection{Recent Systems and the State of the Art}
\label{sec:sota}
Because our central claim concerns comparability rather than rank, we survey
recent RAVDESS systems by \emph{protocol} rather than by score.

Among speaker-dependent designs, Jafarzadeh \emph{et al.}
\cite{jafarzadeh2024ssl} report 97.8--99.0\% for self-supervised encoders with a
lightweight classifier under an explicitly random stratified 80/20 split;
emotion2vec \cite{ma2024emotion2vec} reports 82.43\% under what it calls
``random leave-one-out 10-fold CV'', a phrase that reads as speaker-independent
but denotes a random utterance-level partition; and Ibrahim \emph{et al.}
\cite{ibrahim2026multimodal} report 95.83\% for audio-visual hybrid fusion.

Among speaker-independent designs, CAMuLeNet \cite{goel2024camulenet} reports
82.3\% under ten-fold leave-speaker-out and, valuably, also reports
single-encoder baselines under the identical protocol---58.7\% for HuBERT,
65.4\% for WavLM, 39.6\% for wav2vec~2.0. Luna-Jim{\'e}nez \emph{et al.}
\cite{lunajimenez2022multimodal} report 81.82\% under subject-wise five-fold
cross-validation with published fold assignments. Portal \emph{et al.}
\cite{portal2025benchmarking} re-validate state-of-the-art SER transformers
under strictly speaker-independent conditions and find most single-corpus
results below 40\%. Ibrahim \emph{et al.} \cite{ibrahim2026multimodal} report
both protocols for one system: 95.83\% random against 48.06\%~$\pm$~9.76
leave-one-speaker-out.

Two observations follow. First, the two most informative recent papers are
those that report \emph{both} protocols or several encoders under \emph{one}
protocol; single-number reports are the least usable. Second, the
speaker-independent class spans 44.7 points, wider than the speaker-dependent
class, so protocol annotation alone does not make figures comparable. Our
contribution to this literature is a controlled decomposition: rather than
comparing across papers that differ in many respects, we vary the protocol alone
within one system and attribute the resulting gap to its components.

% =====================================================================
\section{Problem Formulation}
\label{sec:formulation}
% =====================================================================

\subsection{Notation and Task}
Let $\mathcal{A}$ denote the space of speech waveforms and
$\mathcal{Y}=\{1,\dots,C\}$ the set of emotion categories, $C = 8$. A corpus is
a set $\mathcal{D}=\{(a_k, y_k, s_k)\}_{k=1}^{N}$ of utterances $a_k \in
\mathcal{A}$ with labels $y_k \in \mathcal{Y}$ and speaker identifiers $s_k \in
\mathcal{S}$; for RAVDESS, $N = 1440$ and $|\mathcal{S}| = 24$. Two feature maps
extract complementary views of the same signal,
\begin{equation}
  \varphi_p : \mathcal{A} \rightarrow \mathbb{R}^{d_p}, \qquad
  \varphi_h : \mathcal{A} \rightarrow \mathbb{R}^{d_h},
  \label{eq:views}
\end{equation}
with $d_p = 34$ perceptually motivated descriptors and $d_h = 768$ dimensions
from a frozen self-supervised encoder. Write $\mathbf{x}_p = \varphi_p(a)$ and
$\mathbf{x}_h = \varphi_h(a)$.

\subsection{The Unified Representation Problem}
We seek stream encoders $\mathrm{Enc}_p : \mathbb{R}^{d_p}\!\rightarrow\!
\mathbb{R}^{d}$ and $\mathrm{Enc}_h : \mathbb{R}^{d_h}\!\rightarrow\!
\mathbb{R}^{d}$, a fusion operator $\omega : \mathbb{R}^{d} \times
\mathbb{R}^{d} \rightarrow \mathbb{R}^{d}$ drawn from a family $\Omega$, and a
classifier $g_\psi : \mathbb{R}^{d} \rightarrow \Delta^{C-1}$, such that the
\emph{unified representation}
\begin{equation}
  \mathbf{z} = \omega\big(\mathrm{Enc}_p(\mathbf{x}_p),\;
                          \mathrm{Enc}_h(\mathbf{x}_h)\big) \in \mathbb{R}^{d}
  \label{eq:unified}
\end{equation}
minimises the expected risk
\begin{equation}
  (\theta^\star\!, \omega^\star) = \arg\min_{\theta,\, \omega \in \Omega}\;
  \mathbb{E}_{(a,y)\sim P}\big[\ell\big(g_\psi(\mathbf{z}),\, y\big)\big],
  \label{eq:risk}
\end{equation}
subject to the \emph{capacity constraint}
\begin{equation}
  \dim\!\big(\omega(\cdot,\cdot)\big) = d \quad \text{for every } \omega \in
  \Omega .
  \label{eq:capacity}
\end{equation}
Constraint~\eqref{eq:capacity} is what makes an ablation over $\Omega$
interpretable: every operator hands the classifier the same number of
dimensions, so a difference in risk is attributable to the operator and not to
downstream width. We take $\Omega = \{\textsc{concat}, \textsc{gate},
\textsc{xattn}\}$, defined in Section~\ref{sec:fusion-design}.

\subsection{When Can the Psychoacoustic View Contribute?}
\label{sec:whencontrib}
Objective~1 asks whether $\varphi_p$ is worth including at all. Let $A$ be a
random utterance with label $Y$. The relevant quantity is the \emph{unique
information} the perceptual view carries about the label given the
self-supervised view,
\begin{equation}
  U_p \;=\; I\big(\varphi_p(A);\, Y \,\big|\, \varphi_h(A)\big) \;\geq\; 0 ,
  \label{eq:unique}
\end{equation}
so that
\begin{equation}
  I\big(\varphi_p(A), \varphi_h(A);\, Y\big)
  = I\big(\varphi_h(A);\, Y\big) + U_p .
  \label{eq:chain}
\end{equation}
\emph{Complementarity} is exactly the condition $U_p > 0$; \emph{redundancy} is
$U_p = 0$. Two consequences shape the rest of the paper.

First, $U_p > 0$ is necessary but not sufficient for an accuracy gain: a finite
estimator trained on $N = 1440$ utterances may fail to extract information that
is present. A null result on accuracy therefore does not establish redundancy.

Second, $U_p > 0$ implies that, in the population limit, there exist utterances
the perceptual view classifies correctly and the self-supervised view does not.
Writing $E_m$ for the error set of the system trained on view $m$, this
motivates measuring $|E_h \setminus E_p|$ directly rather than inferring
complementarity from a difference of means---a distinction that becomes decisive
once Section~\ref{sec:mde} bounds what the difference of means can resolve.
Algorithm~\ref{alg:complement} implements the measurement.

\subsection{Evaluation Protocol as an Operator}
A \emph{protocol} is a partition operator
\begin{equation}
  \Pi : (\mathcal{D}, s) \;\longmapsto\;
  \big(I_{\mathrm{tr}}, I_{\mathrm{va}}, I_{\mathrm{te}}\big),
  \label{eq:protocol}
\end{equation}
mapping a corpus and a seed to disjoint index sets, together with a rule
naming the partition on which the model is selected. The reported risk is the
protocol-conditioned estimate
\begin{equation}
  \widehat{R}_\Pi(f) \;=\; \frac{1}{|I_{\mathrm{te}}|}
  \sum_{k \in I_{\mathrm{te}}} \mathbb{1}\big[f(a_k) \neq y_k\big].
  \label{eq:risk-hat}
\end{equation}

\begin{definition}[Speaker-disjointness]
$\Pi$ is speaker-disjoint if $\mathcal{S}(I_{\mathrm{tr}}) \cap
\mathcal{S}(I_{\mathrm{te}}) = \emptyset$, where $\mathcal{S}(I) = \{s_k : k \in
I\}$.
\end{definition}

Two effects can now be compared on the same scale. For a fixed architecture
$\mathcal{A}$ and two protocols, the \emph{protocol effect} is
\begin{equation}
  \Delta_\Pi(\mathcal{A}) =
  \big|\widehat{R}_{\Pi_1}(f_{\mathcal{A}}) -
       \widehat{R}_{\Pi_2}(f_{\mathcal{A}})\big| ;
  \label{eq:deltapi}
\end{equation}
for a fixed protocol, the \emph{architecture effect} is
\begin{equation}
  \Delta_\Omega(\Pi) = \max_{\omega, \omega' \in \Omega}
  \big|\widehat{R}_{\Pi}(f_{\omega}) - \widehat{R}_{\Pi}(f_{\omega'})\big| .
  \label{eq:deltaomega}
\end{equation}
The paper's central empirical claim is that $\Delta_\Pi(\mathcal{A}) \gg
\Delta_\Omega(\Pi)$ on this corpus, by a factor of roughly fifteen. Because
$\Delta_\Omega(\Pi)$ is small, the choice of which fixed $\mathcal{A}$ anchors
the ladder in~\eqref{eq:deltapi} shifts the measured protocol effect by at most
$\Delta_\Omega(\Pi)$, which is 2.03 points here.

\subsection{Admissibility and Design Sensitivity}
\label{sec:admissible}
Let $\delta$ be a claimed improvement measured as a mean paired difference over
$n$ folds with standard deviation $\sigma_\Delta$. For a two-sided paired
$t$-test at level $\alpha$ and power $1-\beta$, the \emph{minimum detectable
effect} is
\begin{equation}
  \delta_{\min}(n) \;=\; \sigma_\Delta \cdot d_z^{\star}(n, \alpha, \beta),
  \label{eq:mde}
\end{equation}
where $d_z^{\star}$ solves
$1 - \beta = \Pr\big[\,|T'_{n-1}(d_z\sqrt{n})| > t_{n-1,\,1-\alpha/2}\,\big]$
and $T'$ is the noncentral $t$ distribution.

\begin{definition}[Admissibility]
\label{def:admissible}
A representational or architectural claim of size $\delta$, made on corpus
$\mathcal{D}$ under protocol $\Pi$ with $n$ folds, is \emph{admissible} only if
$|\delta| \geq \delta_{\min}(n)$.
\end{definition}

This is a statement about design, not about an observed $p$-value: post-hoc
power computed from an observed effect is a monotone function of that effect's
$p$-value and carries no additional information \cite{hoenig2001abuse}.
Section~\ref{sec:mde} computes $\delta_{\min}$ for our design and applies
Definition~\ref{def:admissible} to our own results, which it fails.

\subsection{Properties of Pooled-Vector Fusion}
\label{sec:properties}

The three operators in $\Omega$ admit analytical comparison, and two of the
three results below were not anticipated when the framework was designed.

\begin{proposition}[Degeneracy of single-token cross-attention]
\label{prop:degen}
Let $\mathbf{q},\mathbf{k},\mathbf{v} \in \mathbb{R}^{1\times d_k}$ be the
projected query, key and value of a scaled dot-product attention block whose
query and key/value sequences each have length one. Then the attention output
satisfies $\mathbf{A} = \mathbf{v}$ identically, and
$\partial \mathbf{A}/\partial \mathbf{q} = \mathbf{0}$, for any number of heads,
any head dimension and any scaling.
\end{proposition}
\begin{IEEEproof}
The score matrix is $\mathbf{S} = \mathbf{q}\mathbf{k}^{\!\top}/\sqrt{d_k} \in
\mathbb{R}^{1\times 1}$, and
$\mathrm{softmax}(\mathbf{S})_{11} = e^{S_{11}}/e^{S_{11}} = 1$ for every
$S_{11}\in\mathbb{R}$. Hence $\mathbf{A} = 1\cdot\mathbf{v}$, which does not
depend on $\mathbf{q}$; the partial derivative is therefore identically zero.
Multiple heads apply the same argument per head, and any scaling factor cancels
in the ratio.
\end{IEEEproof}

\begin{corollary}
\label{cor:inert}
If $\omega = \textsc{xattn}$ with $K=1$ and the psychoacoustic stream supplies
the query, then $\partial\mathbf{z}/\partial\mathbf{e}_p = \mathbf{0}$: the
psychoacoustic encoder receives no gradient and its parameters remain at
initialisation. The model is \emph{exactly} the unimodal self-supervised model
plus an unused parameter block, and cannot in principle demonstrate
$U_p > 0$.
\end{corollary}

Corollary~\ref{cor:inert} is the reason $K > 1$ is asserted at line~11 of
Algorithm~\ref{alg:unified}; Section~\ref{sec:degeneracy} confirms it
empirically.

\begin{proposition}[Convex-hull confinement of gated fusion]
\label{prop:hull}
For $\mathbf{z} = \mathbf{g}\odot\mathbf{e}_p + (\mathbf{1}-\mathbf{g})
\odot\mathbf{e}_h$ with $\mathbf{g} \in (0,1)^d$,
\begin{equation}
  \min(e_{p,i}, e_{h,i}) < z_i < \max(e_{p,i}, e_{h,i}) \quad \forall i,
\end{equation}
hence $\|\mathbf{z}\|_\infty \leq \max(\|\mathbf{e}_p\|_\infty,
\|\mathbf{e}_h\|_\infty)$. No such bound holds for
$\mathbf{z} = \mathbf{W}_c[\mathbf{e}_p;\mathbf{e}_h]$. Consequently neither
operator's function class contains the other.
\end{proposition}
\begin{IEEEproof}
Each $z_i$ is a strict convex combination of $e_{p,i}$ and $e_{h,i}$ because
$g_i \in (0,1)$, giving the bound. For the converse, take $\mathbf{W}_c$ with
$\mathbf{z} = 2(\mathbf{e}_p + \mathbf{e}_h)$, whose $\infty$-norm exceeds the
bound whenever $\mathbf{e}_p,\mathbf{e}_h \neq \mathbf{0}$ share a sign in some
coordinate; and note that $\mathbf{g}$ depends nonlinearly on
$[\mathbf{e}_p;\mathbf{e}_h]$, so the gated map is not linear and is not
realisable by any fixed $\mathbf{W}_c$.
\end{IEEEproof}

Proposition~\ref{prop:hull} explains the empirical near-tie of
Section~\ref{sec:fusion} without appealing to noise alone: the two operators
span overlapping but distinct function classes, and the emotion-relevant signal
in a pair of pooled vectors evidently lies in the intersection. It also
identifies what the gate buys that concatenation does not---an explicit
per-coordinate attribution $\mathbf{g}$ between the perceptual and semantic
views, which Section~\ref{sec:modality} analyses.

\begin{proposition}[Cost ordering]
\label{prop:cost}
With encoder width $d$, token count $K$, and unshared expansions and
projections in both attention directions, the trainable fusion parameter counts
are $2d^2 + d$ for both \textsc{concat} and \textsc{gate}, and
$4(Kd^2 + Kd) + 2(4d^2+4d) + 4d + (2d^2+d)$ for \textsc{xattn}. At $d = 128$,
$K = 4$ this gives $32{,}896$ against $429{,}696$, and total trainable counts of
$187{,}656$ against $584{,}456$, a ratio of $3.11\times$; the corresponding
multiply--accumulate ratio per utterance is $5.28\times$
(Table~\ref{tab:complexity}).
\end{proposition}

% =====================================================================
\begin{figure*}[!t]
\centering
\paperfig{figures/fig1_framework_ieee}{\wfigw}{1.053}
\caption{The PAER-LLM framework. Two feature views of one audio signal---34
psychoacoustic descriptors and a 768-dimensional frozen HuBERT-base
embedding---are independently encoded to 128 dimensions and combined by a
learned per-dimension gate into a single unified representation $\mathbf{z} \in
\mathbb{R}^{128}$. Two alternative operators, concatenation and token
cross-attention, are evaluated and not adopted. The resulting posterior over
eight categories conditions an instruction-tuned language model; only the
posterior is passed, and learned activations are not. Layers~4 and~5 are drawn
dashed, below the divider, because they are proposed and implemented but not
evaluated in this work. Edge annotations give tensor shapes.}
\label{fig:framework}
\end{figure*}

\emph{Reading Fig.~\ref{fig:framework}.} The figure is read top to bottom as a
pipeline and left to right within each layer. The left tab names the layer; the
right panel states what that layer is responsible for delivering. Layer~1 ends
in a verified corpus with duplication factor $r=1$, which is the precondition
for every number in this paper. Layer~2 shows the two views side by side with
their dimensionalities, and the ``$+$'' between them denotes that they are
carried forward jointly rather than combined at this stage. Layer~3 is the
paper's object of study: the two encoders map $34$ and $768$ dimensions to a
common $128$, the gate combines them, and the classifier emits an eight-way
posterior; the italic line beneath records the two operators evaluated and not
adopted. The dashed divider below Layer~3 is the boundary of the experimental
evidence---everything above it is measured, everything below it is specified and
implemented but not evaluated, and the dashed outlines of Layers~4 and~5 carry
that distinction visually rather than relying on the reader to recall it from
the text.

% =====================================================================
\section{The PAER-LLM Framework}
\label{sec:framework}
% =====================================================================
Fig.~\ref{fig:framework} summarises the whole; we describe each layer in turn.

\subsection{Data Acquisition Layer}
\label{sec:data}
The corpus is RAVDESS \cite{livingstone2018ravdess}, restricted to speech: 1440
utterances from 24 actors (12 female, 12 male) across eight categories. All
categories except \emph{neutral} contain 192 utterances; \emph{neutral} contains
96, as it is recorded at a single intensity. The corpus also contains 1012 song
recordings, which we exclude. Audio is resampled to 22{,}050~Hz for
psychoacoustic analysis and to 16{,}000~Hz for the self-supervised encoder,
matching its pretraining condition.

\emph{Corpus integrity.} We recommend that this layer include an explicit
integrity check. In our case the distributed archive had been unpacked twice
into nested directories, so a recursive file walk returned 2880 paths for 1440
distinct recordings and every derived feature table contained each utterance
exactly twice. Because the duplicate rows are byte-identical, no summary
statistic reveals the problem: class balance, feature distributions and file
counts all appear plausible. The defect surfaces only when the number of unique
filenames is compared against the number of rows. We report it because it is
silent, easy to introduce, and materially inflates every downstream result. The
released harness performs the check and refuses to proceed on a mismatch
(Algorithm~\ref{alg:integrity}), returning the duplication factor $r$ so that
the contaminated condition can be reconstructed deliberately for protocol~P1.

\begin{algorithm}[!t]
\caption{Corpus integrity verification}
\label{alg:integrity}
\begin{algorithmic}[1]
\Require dataset root $D$; expected corpus size $N_{\mathrm{exp}}$
\Ensure unique utterance set $U$; duplication factor $r$
\State $P \gets \textsc{Sort}(\text{all \texttt{.wav} paths under } D)$
\State $U \gets \emptyset$; \; $\mathcal{S} \gets \emptyset$
\ForAll{$p \in P$}
  \State $b \gets \textsc{Basename}(p)$
  \If{$b \notin \mathcal{S}$}
    \State $\mathcal{S} \gets \mathcal{S} \cup \{b\}$; \;
           $U \gets U \cup \{p\}$
  \EndIf
\EndFor
\State $r \gets |P| / |U|$
\If{$r > 1$}
  \State \textbf{abort} unless duplication is explicitly requested
  \Statex \Comment{$r>1$: exact copies of test items enter training}
\EndIf
\If{$|U| \neq N_{\mathrm{exp}}$} \State \textbf{warn} \EndIf
\State \Return $U,\; r$
\end{algorithmic}
\end{algorithm}

\subsection{Feature Representation Layer}
\label{sec:features}

\emph{Psychoacoustic stream.} We extract 34 perceptually motivated descriptors
per utterance. \emph{Pitch} (5): mean, maximum, minimum, standard deviation and
range of voiced $F_0$, from Praat's autocorrelation tracker over 75--600~Hz. No
utterance is wholly unvoiced, so the unvoiced fallback is never exercised.
\emph{Intensity} (5): the same five statistics of frame-level RMS energy. We
call these \emph{intensity} and not \emph{loudness}: loudness in the
psychoacoustic sense \cite{iso532} involves critical-band excitation, masking
and a compressive sone scale, none of which RMS energy models. \emph{Bark-band
energies} (24): mean energy in each of 24 bands on the critical-band-rate scale
of Zwicker and Terhardt \cite{zwicker1980analytical},
\begin{equation}
  z(f) = 13\arctan(0.00076 f)
       + 3.5\arctan\!\left(\left(\frac{f}{7500}\right)^{\!2}\right),
  \label{eq:bark}
\end{equation}
where $f$ is in Hz and $z$ in Bark. Note that the argument of the second
arctangent is squared, not the arctangent itself. Triangular filters are
area-normalised so band energy is not biased by the increasing bandwidth of
higher bands.

One deviation requires comment. The 24 Zwicker bands extend to 15.5~kHz
($z = 24.0$), above the 11.025~kHz Nyquist frequency of our 22.05~kHz analysis
rate. Rather than retain bands that would be identically zero, we place 24
equal-width bands over $[0, z(f_{\mathrm{Nyq}})] = [0, 22.86]$~Bark, so all lie
below Nyquist by construction. The cost is that the upper $1.14$~Bark of the
audible range is unrepresented and the bands are not the canonical Zwicker
bands. Since RAVDESS is distributed at 48~kHz, analysing at $\geq 32$~kHz would
remove the constraint entirely; we regard this as a defect of the present
configuration rather than a design choice to defend.
Algorithm~\ref{alg:bark} gives the construction.

\begin{algorithm}[!t]
\caption{Nyquist-adaptive Bark filter bank}
\label{alg:bark}
\begin{algorithmic}[1]
\Require sample rate $f_s$; FFT size $n$; band count $B = 24$
\Ensure filter bank $\Phi \in \mathbb{R}^{B \times (1 + n/2)}$
\State $\mathbf{f} \gets$ FFT bin centres on $[0, f_s/2]$; \;
       $\mathbf{z} \gets z(\mathbf{f})$ \Comment{Eq.~\eqref{eq:bark}}
\State $z_{\max} \gets z(f_s/2)$
  \Comment{adaptive; canonical scheme reaches $24.0$ at 15.5~kHz}
\State $\mathbf{e} \gets \textsc{Linspace}(0, z_{\max}, B+2)$
\For{$i = 1$ \textbf{to} $B$}
  \State $(\ell, c, u) \gets (e_i, e_{i+1}, e_{i+2})$
  \State $\Phi_{i,:} \gets \max\!\big(0, \min(\tfrac{\mathbf{z}-\ell}{c-\ell},
         \tfrac{u-\mathbf{z}}{u-c})\big)$
  \State $\Phi_{i,:} \gets \Phi_{i,:} / \textstyle\sum_j \Phi_{i,j}$
    \Comment{unit area}
\EndFor
\State \textbf{verify} $\mathrm{Var}_k(\Phi_{i,:}\mathbf{S}_k) > 0\;\forall i$
\State \Return $\Phi$
\end{algorithmic}
\end{algorithm}

\emph{Self-supervised stream.} HuBERT-base \cite{hsu2021hubert}, pretrained on
LibriSpeech-960 and held frozen. The final hidden state is mean-pooled over time
to give a 768-dimensional utterance embedding.

\subsection{Unified Representation Layer}
\label{sec:fusion-design}
Each stream is encoded by an independent network
$\mathbf{e}_m = \mathrm{Enc}_m(\mathbf{x}_m) \in \mathbb{R}^{d}$, $m\in\{p,h\}$,
with $d = 128$ and the stack
$\mathrm{Linear}(d_m\!\to\! d) \rightarrow \mathrm{LayerNorm} \rightarrow
\mathrm{GELU} \rightarrow \mathrm{Dropout}(0.3) \rightarrow
\mathrm{Linear}(d\!\to\! d) \rightarrow \mathrm{LayerNorm} \rightarrow
\mathrm{GELU}$, using post-normalisation, PyTorch default initialisation and no
residual connection \cite{ba2016layernorm,hendrycks2016gelu,srivastava2014dropout}.
Encoding both streams to a common width before fusion is what satisfies the
capacity constraint~\eqref{eq:capacity}. We compare three operators.

\emph{Concatenation.} $\mathbf{z} = \mathbf{W}_c[\mathbf{e}_p;\mathbf{e}_h]$
with $\mathbf{W}_c \in \mathbb{R}^{2d\times d}$.

\emph{Gated fusion (adopted).}
\begin{equation}
  \mathbf{g} = \mathrm{sigm}\!\big(\mathbf{W}_g[\mathbf{e}_p;\mathbf{e}_h]
    + \mathbf{b}\big), \quad
  \mathbf{z} = \mathbf{g}\odot\mathbf{e}_p + (\mathbf{1}-\mathbf{g})
  \odot\mathbf{e}_h .
  \label{eq:gate}
\end{equation}
Each coordinate of $\mathbf{g}$ is the weight assigned to the perceptual view at
that coordinate, so $\mathbf{g}$ is a directly readable per-dimension allocation
between the two representations---a property neither alternative has
(Proposition~\ref{prop:hull}).

\emph{Token cross-attention.} Each pooled vector is expanded into $K = 4$
learned tokens,
\begin{equation*}
  \mathbf{T}_m = \mathrm{reshape}(\mathbf{W}_m\mathbf{e}_m)
  \in \mathbb{R}^{K\times d}, \qquad m \in \{p, h\},
\end{equation*}
and the streams attend bidirectionally,
\begin{equation}
  \mathbf{A}_{p\rightarrow h} = \mathrm{MHA}(\mathbf{T}_p,\mathbf{T}_h,
  \mathbf{T}_h), \quad
  \mathbf{A}_{h\rightarrow p} = \mathrm{MHA}(\mathbf{T}_h,\mathbf{T}_p,
  \mathbf{T}_p),
  \label{eq:xattn}
\end{equation}
with four heads of dimension 32 and attention dropout 0.3. The expansions and
the attention projections are unshared between directions, so the operator
instantiates four expansion matrices and two complete attention blocks; a
post-norm residual, mean pooling over tokens, concatenation and a projection
$\mathbb{R}^{2d}\!\to\!\mathbb{R}^{d}$ complete it. This accounts exactly for the
$584{,}456$ trainable parameters of Proposition~\ref{prop:cost}. The expansion
to $K > 1$ is essential by Corollary~\ref{cor:inert}.

All three operators emit exactly $d = 128$ dimensions, so the head that follows
has identical capacity ($17{,}544$ parameters) in every condition.
$\mathbf{z}$ is layer-normalised and passed to
$\mathrm{Dropout} \rightarrow \mathrm{Linear}(d\!\to\!d) \rightarrow
\mathrm{GELU} \rightarrow \mathrm{Dropout} \rightarrow
\mathrm{Linear}(d\!\to\!8)$. The softmax posterior is retained in full.

\begin{algorithm}[!t]
\caption{Unified psychoacoustic--semantic representation}
\label{alg:unified}
\begin{algorithmic}[1]
\Require $\mathbf{x}_p \in \mathbb{R}^{34}$, $\mathbf{x}_h \in \mathbb{R}^{768}$;
         scalers $(S_p,S_h)$; width $d$; operator $\omega$
\Ensure representation $\mathbf{z}$; posterior $\hat{\mathbf{p}}$; gate
        $\mathbf{g}$
\State $\mathbf{x}_p \gets S_p(\mathbf{x}_p)$; \;
       $\mathbf{x}_h \gets S_h(\mathbf{x}_h)$
       \Comment{fitted on training rows only}
\State $\mathbf{e}_p \gets \mathrm{Enc}_p(\mathbf{x}_p)$; \;
       $\mathbf{e}_h \gets \mathrm{Enc}_h(\mathbf{x}_h)$
       \Comment{both to $\mathbb{R}^{d}$, Eq.~\eqref{eq:capacity}}
\State $\mathbf{g} \gets \bot$
\If{$\omega = \textsc{gate}$}
  \State $\mathbf{g} \gets \mathrm{sigm}(\mathbf{W}_g[\mathbf{e}_p;
         \mathbf{e}_h]+\mathbf{b})$
  \State $\mathbf{z} \gets \mathbf{g}\odot\mathbf{e}_p +
         (\mathbf{1}-\mathbf{g})\odot\mathbf{e}_h$
         \Comment{Eq.~\eqref{eq:gate}}
\ElsIf{$\omega = \textsc{concat}$}
  \State $\mathbf{z} \gets \mathbf{W}_c[\mathbf{e}_p;\mathbf{e}_h]$
\Else
  \State $\mathbf{T}_p \gets \mathrm{reshape}(\mathbf{W}_p\mathbf{e}_p)$; \;
         $\mathbf{T}_h \gets \mathrm{reshape}(\mathbf{W}_h\mathbf{e}_h)$
  \State \textbf{assert} $K > 1$
    \Comment{$K{=}1$ is gradient-inert, Cor.~\ref{cor:inert}}
  \State $\mathbf{A}_{p\rightarrow h},\mathbf{A}_{h\rightarrow p} \gets$
         Eq.~\eqref{eq:xattn}
  \State $\mathbf{z} \gets \mathbf{W}_o[\overline{\mathbf{A}}_{p\rightarrow h};
         \overline{\mathbf{A}}_{h\rightarrow p}]$
\EndIf
\State $\mathbf{z} \gets \mathrm{LayerNorm}(\mathbf{z})$;\;
       $\hat{\mathbf{p}} \gets \mathrm{softmax}(\mathrm{Head}(\mathbf{z}))$
\State \Return $\mathbf{z},\, \hat{\mathbf{p}},\, \mathbf{g}$
\end{algorithmic}
\end{algorithm}

\subsection{Language Intelligence Layer}
\label{sec:llm}
The predicted distribution conditions an instruction-tuned language model,
Phi-3-mini-4k-instruct \cite{abdin2024phi3}. Rather than passing the top-1
label, the prompt encodes the three highest-probability categories with their
posterior mass and instructs the model to hedge when the top two lie within a
margin $\tau$ (Algorithm~\ref{alg:prompt}). The intent is that an ambiguous
vocal signal should produce a tentative response rather than a confident
misattribution \cite{tian2023calibration,wen2025abstention}. We do \emph{not}
pass the learned representation $\mathbf{z}$ to the language model: a model not
trained alongside that network has no basis for interpreting its activations.

We state plainly that this layer is unevaluated here---no human study, no
automatic metric, no ablation against top-1-conditioned prompting---and
$\tau$ is asserted at $0.15$ rather than tuned. Algorithm~\ref{alg:prompt}
therefore treats $\tau$ as a free parameter, to be set from the
margin--accuracy curve of Algorithm~\ref{alg:complement}.

\begin{algorithm}[!t]
\caption{Posterior-conditioned prompt construction}
\label{alg:prompt}
\begin{algorithmic}[1]
\Require posterior $\hat{\mathbf{p}}$; query $q$; context $h$; margin $\tau$;
         model $\mathcal{M}$
\Ensure emotion-aware response $y$
\State $(c_1,c_2,c_3) \gets$ indices of the three largest entries of
       $\hat{\mathbf{p}}$
\State $\mu \gets \hat{p}_{c_1} - \hat{p}_{c_2}$ \Comment{top-2 margin}
\If{$\mu < \tau$}
  \State $s \gets$ ``evidence is ambiguous between $c_1$ and $c_2$;
         acknowledge the uncertainty''
\Else
  \State $s \gets$ ``evidence indicates $c_1$; respond to that state''
\EndIf
\State $\pi \gets \textsc{Compose}(\{(c_i,\hat{p}_{c_i})\}_{i=1}^{3}, s, h, q)$
  \Comment{mass stated verbatim}
\State $y \gets \mathcal{M}(\pi)$; \;
       \textbf{record} $(\hat{\mathbf{p}}, \mu)$ for calibration
\State \Return $y$
\end{algorithmic}
\end{algorithm}

\subsection{Application Layer}
The pipeline accepts an utterance, returns a predicted category with its
posterior, and generates a corresponding response. A demonstration interface is
included in the released code. As with Layer~4, no user study supports this
layer and we make no claim for it.

% =====================================================================
\section{Research Methodology}
\label{sec:methodology}
% =====================================================================

\subsection{Methodological Stance}
The methodology follows from Section~\ref{sec:formulation}. Because
$\Delta_\Pi(\mathcal{A})$ and $\Delta_\Omega(\Pi)$ are defined on the same
scale, they can be measured in one experimental design by holding one factor
fixed while varying the other. We therefore run two orthogonal sweeps on the
same code path, the same cached features and the same seeds: a
\emph{protocol sweep} with the architecture fixed, and an \emph{architecture
sweep} with the protocol fixed at its most defensible setting. A third sweep
varies the input streams to address $U_p$ directly. Nothing else changes between
runs; the ablations are therefore controlled rather than comparative.

\subsection{Design Decisions and the Alternatives Considered}
Table~\ref{tab:methodology} records each substantive methodological choice, the
alternatives available, and the reason for the decision. Two of the decisions
work against the apparent strength of our results and are marked accordingly;
we record them because a methodology section that lists only favourable choices
is not a methodology section.

\begin{table*}[!t]
\centering
\small
\caption{Research methodology: design decisions, alternatives considered, and
rationale. Rows marked $\ominus$ are choices that reduce our reported
performance or weaken our own claim, and were adopted anyway.}
\label{tab:methodology}
\renewcommand{\arraystretch}{1.15}
\begin{tabular}{@{}L{0.15\textwidth}L{0.24\textwidth}L{0.19\textwidth}L{0.33\textwidth}@{}}
\toprule
\textbf{Decision} & \textbf{Alternatives considered} & \textbf{Adopted} &
\textbf{Rationale} \\
\midrule
Encoder treatment $\ominus$ & fine-tuned encoder; weighted layer combination;
frozen final layer & frozen final layer, mean-pooled &
Isolates the fusion operator and the psychoacoustic stream from encoder
adaptation, which is known to dominate \cite{chen2023exploring}. Costs us
absolute accuracy; our figures are a floor. \\
Perceptual descriptors & eGeMAPS-88; ISO 532-1 loudness/sharpness/roughness;
cochleagram front-end & 34 descriptors (5 $F_0$, 5 RMS, 24 Bark) &
Small, interpretable and directly tied to critical-band theory. Benchmarked
against eGeMAPS on identical folds rather than assumed adequate
(Table~\ref{tab:modality}). \\
Fusion family $\Omega$ & attention-only; concatenation-only; late (decision)
fusion & all three, at fixed output width &
Constraint~\eqref{eq:capacity} makes the ablation attributable to the operator.
Late fusion was excluded because it changes head capacity. \\
Primary protocol & random split; single actor-disjoint split; LOSO &
actor-disjoint 5-fold &
Speaker-disjoint by Definition~1, with dispersion; LOSO would be better
(Section~\ref{sec:mde}) but was not affordable in the present grid. \\
Protocol ladder anchor & re-run all rungs per operator & one fixed operator
(token cross-attention) &
$\Delta_\Pi$ is defined at fixed $\mathcal{A}$; by~\eqref{eq:deltaomega} the
anchor shifts the result by at most 2.03 points. \\
Model selection & test partition; validation partition &
validation, with test selection reproduced only to measure its optimism &
Test selection is the leakage mode most often undisclosed; P2t/P3t exist to
quantify it and are never used for a reported result. \\
Primary metric & accuracy; UAR; macro-$F_1$ & macro-$F_1$ &
Fixed in advance because it is also the model-selection criterion, and
\emph{neutral} has half the support. \\
Multiplicity control $\ominus$ & none; Holm across all six tests &
Holm within each metric (3 tests) &
Accuracy and macro-$F_1$ on identical folds are strongly dependent; six tests
is neither exact nor conservative. This is stricter than no correction. \\
Claim threshold $\ominus$ & report $p<0.05$; report point estimates &
admissibility, Definition~\ref{def:admissible} &
Applied to our own results, which fail it (Section~\ref{sec:mde}). \\
Complementarity test & difference of means & error sets, McNemar, per-class
$\Delta$recall, gate distribution &
$U_p > 0$ can hold below $\delta_{\min}$; means cannot detect it
(Section~\ref{sec:whencontrib}). \\
\bottomrule
\end{tabular}
\end{table*}

\subsection{Benchmarking Methodology}
\label{sec:benchmark-method}
Benchmarking in this paper is \emph{protocol-stratified} and \emph{in-harness},
for reasons that follow from~\eqref{eq:deltapi}.

\emph{Protocol-stratified external benchmarking.} Published figures are grouped
by protocol class and never compared across classes
(Table~\ref{tab:matrix}). Within a class they remain only loosely
comparable, because encoder adaptation, fold count and model-selection practice
still differ; we therefore additionally record whether each source reports any
measure of dispersion, since a point estimate without dispersion cannot support
a ranking at the margins the literature disputes.

\emph{In-harness baselines.} Any comparison we make on our own behalf is against
a system we ran ourselves, on our own folds, in our own code path. Three
baselines serve this purpose: chance ($12.50\%$); an MFCC-40 plus random-forest
classifier, representing the pre-deep-learning approach; and frozen HuBERT-base
alone, which is the ablation of our own psychoacoustic stream and simultaneously
the strongest single-encoder reference. The HuBERT-only condition is the one
that matters for Objective~1, because it is the term $I(\varphi_h(A);Y)$
in~\eqref{eq:chain} made operational.

\emph{What we do not do.} We do not report a best-of-$n$ figure, do not select
the fusion operator on the test partition, and do not compare our
speaker-independent numbers against speaker-dependent published numbers. The
last of these would improve our apparent standing considerably.

\subsection{Controlled Decomposition Procedure}
Algorithm~\ref{alg:decompose} states the protocol sweep. Three configuration
switches vary; everything else is held fixed, so each increment is attributable
to a single evaluation-design choice.

\begin{algorithm}[!t]
\caption{Controlled decomposition of the protocol effect}
\label{alg:decompose}
\begin{algorithmic}[1]
\Require features $(X_p,X_h)$, labels $y$, metadata $M$; seeds $\mathcal{S}$;
         fixed architecture $\mathcal{A}$
\Ensure accuracy per rung and attributable increments
\State $\mathcal{R} \gets [\,$P1$, $P2$, $P2t$, $P3$, $P3t$, $P4$\,]$
       \Comment{Table~\ref{tab:protocols}}
\ForAll{rungs $\rho \in \mathcal{R}$}
  \State \textbf{if} $\rho.\mathrm{dup}$ \textbf{then}
         $(X,y,M) \gets \textsc{Replicate}(X,y,M;r)$
         \Comment{restore contamination}
  \ForAll{$s \in \mathcal{S}$}
    \State $(I_{\mathrm{tr}},I_{\mathrm{va}},I_{\mathrm{te}}) \gets
           \Pi_{\rho}(\mathcal{D}, s)$ \Comment{Eq.~\eqref{eq:protocol}}
    \State $I_{\mathrm{sel}} \gets I_{\mathrm{te}}$ \textbf{if}
           $\rho.\mathrm{sel}=\mathbf{test}$ \textbf{else} $I_{\mathrm{va}}$
    \State $\theta^\star \gets \textsc{FitFold}(\mathcal{A},I_{\mathrm{tr}},
           I_{\mathrm{va}},I_{\mathrm{sel}},s)$
    \State record accuracy, macro-$F_1$, UAR on $I_{\mathrm{te}}$
  \EndFor
  \State $\mu_\rho \gets$ mean over seeds (and folds for P4)
\EndFor
\State $\Delta_{\mathrm{dup}} \gets \mu_{\mathrm{P1}}-\mu_{\mathrm{P2}}$
  \Comment{local defect, not a field practice}
\State $\Delta_{\mathrm{spk}} \gets \mu_{\mathrm{P2}}-\mu_{\mathrm{P3}}$
  \Comment{$\dagger$}
\State $\Delta_{\mathrm{sel}} \gets \mu_{\mathrm{P3t}}-\mu_{\mathrm{P3}}$
  \Comment{$\dagger$}
\State $\Delta_{\mathrm{split}} \gets \mu_{\mathrm{P3}}-\mu_{\mathrm{P4}}$
  \Comment{$\dagger$}
\State \Return $\{\mu_\rho\},\{\Delta_\bullet\}$
\end{algorithmic}
\end{algorithm}

\subsection{Assessing the Unified Representation}
\label{sec:complement-method}
Aggregate accuracy is a weak instrument for the question of
Section~\ref{sec:whencontrib}: two views can satisfy $U_p > 0$ while their
combination moves the mean by less than $\delta_{\min}$.
Algorithm~\ref{alg:complement} therefore measures complementarity on the error
sets, and computes the calibration diagnostics Layer~4 requires.

\begin{algorithm}[!t]
\caption{Stream complementarity and posterior calibration}
\label{alg:complement}
\begin{algorithmic}[1]
\Require pooled CV predictions $\{(\hat{\mathbf{p}}^{(m)}_k,y_k)\}$ for
         $m\in\{p,h,\mathrm{fused}\}$; gates $\{\mathbf{g}_k\}$; bins $B$
\Ensure complementarity and calibration report
\State $E_m \gets \{k : \arg\max\hat{\mathbf{p}}^{(m)}_k \neq y_k\}$
\State $J \gets |E_p \cap E_h| / |E_p \cup E_h|$
  \Comment{$J\!\to\!1$: the views fail together}
\State \textbf{report} $|E_h\setminus E_p|$, $|E_p\setminus E_h|$
  \Comment{operationalises $U_p>0$, Sec.~\ref{sec:whencontrib}}
\State $m^\star \gets \arg\min_{m\in\{p,h\}}|E_m|$
\State $(b,c) \gets (|E_{m^\star}\!\setminus\! E_{\mathrm{fused}}|,\,
       |E_{\mathrm{fused}}\!\setminus\! E_{m^\star}|)$
\State $p_{\mathrm{McN}} \gets$ exact binomial test on $(b,c)$
\State \textbf{report} per-class $\Delta$recall, fused $-\, m^\star$
\State \textbf{report} mean and spread of $\mathbf{g}$ per dimension, per class
\State $\mathrm{ECE} \gets \sum_{b}\frac{|\mathcal{B}_b|}{N}
       |\mathrm{acc}(\mathcal{B}_b)-\mathrm{conf}(\mathcal{B}_b)|$
\State plot accuracy against margin $\mu$; set $\tau^\star$
  \Comment{fixes $\tau$ in Alg.~\ref{alg:prompt}}
\State \textbf{repeat} ECE under P1 and P4
  \Comment{does contamination change calibration?}
\State \Return $\{J,b,c,p_{\mathrm{McN}},\mathrm{ECE}_{\mathrm{P1}},
       \mathrm{ECE}_{\mathrm{P4}},\tau^\star\}$
\end{algorithmic}
\end{algorithm}

\subsection{Metrics and Statistical Treatment}
\label{sec:stats}
We report accuracy, macro-averaged $F_1$ and unweighted average recall (UAR).
Because \emph{neutral} has half the support of the other classes, accuracy alone
is misleading. For single-split results we report percentile bootstrap 95\%
confidence intervals over 2000 resamples of the test partition. For
cross-validated results we report the mean, the standard deviation across folds
(computed with $n-1$ in the denominator), and the 95\% confidence interval of
the mean; the last is the inferentially meaningful quantity and the three are
labelled distinctly throughout \cite{agarwal2021precipice}.

Paired comparisons between operators use identical folds and report the mean
paired difference, its standard deviation, a 95\% confidence interval, a paired
two-sided $t$-test, and Holm--Bonferroni adjusted $p$-values \emph{within} each
metric across its three pairwise comparisons \cite{demsar2006statistical,
dietterich1998tests}. With five folds the exact Wilcoxon signed-rank test cannot
attain $p < 0.0625$ and is omitted \cite{card2020power}. Design sensitivity
follows~\eqref{eq:mde}.

% =====================================================================
\section{Experimental Setup}
\label{sec:setup}
% =====================================================================

\subsection{Evaluation Protocols}
Six protocols differ only in partitioning and in which partition selects the
model (Table~\ref{tab:protocols}). \emph{Actor} refers to a RAVDESS performer;
\emph{speaker-disjoint} means no actor's utterances appear in more than one
partition. Under all protocols, feature standardisation is fitted on the
training partition alone.

\begin{table}[!t]
\centering\small
\caption{The six evaluation protocols. Architecture, features, optimiser,
schedule and seeds are identical across rows. P2t and P3t reproduce undisclosed
test-set model selection so that its optimism can be measured; they are never
used for a reported result.}
\label{tab:protocols}
\begin{tabular}{@{}llll@{}}
\toprule
\textbf{Protocol} & \textbf{Duplicates} & \textbf{Partition} &
\textbf{Selection} \\
\midrule
P1  & present (2880) & random 80/20   & validation \\
P2  & removed (1440) & random 80/20   & validation \\
P2t & removed        & random 80/20   & \textbf{test} \\
P3  & removed        & actor-disjoint & validation \\
P3t & removed        & actor-disjoint & \textbf{test} \\
P4  & removed        & actor-disjoint & validation, 5-fold \\
\bottomrule
\end{tabular}
\end{table}

\begin{table}[!t]
\centering\small
\caption{Published actor assignments (seed 42). RAVDESS numbers actors so that
odd identifiers are male and even are female.}
\label{tab:folds}
\begin{tabular}{@{}ll@{}}
\toprule
\textbf{Partition} & \textbf{Actors} \\
\midrule
P3 train      & 1, 3, 4, 6, 7, 8, 10, 12, 13, 15, \\
              & 17, 18, 19, 21, 22, 24 \\
P3 validation & 5, 9, 14, 20 \\
P3 test       & 2, 11, 16, 23 \\
\midrule
P4 fold 1 (test) & 3, 8, 11, 23, 24 \\
P4 fold 2 (test) & 6, 7, 14, 18, 19 \\
P4 fold 3 (test) & 1, 9, 16, 21, 22 \\
P4 fold 4 (test) & 2, 10, 12, 13, 15 \\
P4 fold 5 (test) & 4, 5, 17, 20 \\
\bottomrule
\end{tabular}
\end{table}

\begin{algorithm}[!t]
\caption{Gender-balanced speaker-disjoint partitioning}
\label{alg:split}
\begin{algorithmic}[1]
\Require metadata $M$ (\textit{actor}, \textit{gender}); fractions
         $\tau_{\mathrm{te}},\nu$; seed $s$
\Ensure disjoint $I_{\mathrm{tr}},I_{\mathrm{va}},I_{\mathrm{te}}$
\State $A_{\mathrm{tr}},A_{\mathrm{va}},A_{\mathrm{te}} \gets
       \emptyset,\emptyset,\emptyset$
\ForAll{genders $g$ in $M$}
  \State $A_g \gets \textsc{Shuffle}(\text{actors of gender } g;\, s)$
  \State $n_{\mathrm{te}} \gets \max(1,\lfloor |A_g|\tau_{\mathrm{te}}\rceil)$;\;
         $n_{\mathrm{va}} \gets \max(1,\lfloor (|A_g|-n_{\mathrm{te}})\nu\rceil)$
  \State append the first $n_{\mathrm{te}}$ to $A_{\mathrm{te}}$, the next
         $n_{\mathrm{va}}$ to $A_{\mathrm{va}}$, the rest to $A_{\mathrm{tr}}$
\EndFor
\State \textbf{assert} the three actor sets are pairwise disjoint
  \Comment{overlap fails loudly, not silently}
\State $I_\bullet \gets \{k : M_k.\textit{actor} \in A_\bullet\}$
\State \textbf{publish} $A_{\mathrm{tr}},A_{\mathrm{va}},A_{\mathrm{te}}$
  \Comment{Table~\ref{tab:folds}}
\State \Return $I_{\mathrm{tr}},I_{\mathrm{va}},I_{\mathrm{te}}$
\end{algorithmic}
\end{algorithm}

\begin{algorithm}[!t]
\caption{Training one fold without leakage}
\label{alg:fit}
\begin{algorithmic}[1]
\Require $X_p,X_h,y$; $I_{\mathrm{tr}},I_{\mathrm{va}}$; selection set
         $I_{\mathrm{sel}}$ (default $I_{\mathrm{va}}$); epochs $E$;
         patience $\pi$; seed $s$
\Ensure $\theta^\star$; scalers $(S_p,S_h)$
\State $\textsc{SetSeed}(s)$;\;
       $S_m \gets \textsc{Fit}(X_m[I_{\mathrm{tr}}])$ for $m\in\{p,h\}$
       \Comment{training rows only}
\State $\mathbf{w} \gets$ inverse class frequency of $y[I_{\mathrm{tr}}]$
\State $b^\star \gets -\infty$;\; $c \gets \pi$
\For{$e = 1$ \textbf{to} $E$}
  \State train one epoch (Alg.~\ref{alg:unified} forward, weighted
         cross-entropy, gradient clipping)
  \State $m \gets \textsc{MacroF1}$ on $I_{\mathrm{sel}}$
  \If{$m > b^\star$} \State $b^\star \gets m$;\; $\theta^\star \gets \theta$;\;
      $c \gets \pi$
  \Else \State $c \gets c-1$; \textbf{if} $c=0$ \textbf{then break} \EndIf
\EndFor
\State \Return $\theta^\star,(S_p,S_h)$
  \Comment{test partition read exactly once, afterwards}
\end{algorithmic}
\end{algorithm}

\subsection{Implementation}
Table~\ref{tab:hyper} gives the configuration. Optimisation uses AdamW
\cite{loshchilov2019adamw} with cosine annealing \cite{loshchilov2017sgdr} and
label smoothing \cite{szegedy2016rethinking}; cross-entropy is weighted by
inverse class frequency to compensate for the half-sized \emph{neutral} class.

\begin{table}[!t]
\centering\small
\caption{Hyperparameters and signal-processing configuration. HuBERT-base is
frozen; parameter counts exclude it.}
\label{tab:hyper}
\begin{tabular}{@{}ll@{}}
\toprule
Optimiser / LR / weight decay & AdamW / $10^{-3}$ / $10^{-2}$ \\
Schedule / batch / clipping   & cosine annealing / 32 / 5.0 \\
Dropout / label smoothing     & 0.3 / 0.05 \\
Epochs / patience / monitor   & 100 / 15 / val.\ macro-$F_1$ \\
Class weighting               & inverse frequency \\
Width $d$ / tokens $K$ / heads & 128 / 4 / 4 \\
Seeds                         & \TORUN{5 seeds} \\
\midrule
Psychoacoustic sample rate    & 22{,}050 Hz \\
HuBERT sample rate            & 16{,}000 Hz \\
FFT / hop / window            & 2048 / 512 / Hann \\
$F_0$ estimator / range       & Praat autocorr.\ / 75--600 Hz \\
Bark bands / normalisation    & 24 / unit area \\
\midrule
Hardware / software           & \TOADD{GPU, CUDA, versions} \\
\bottomrule
\end{tabular}
\end{table}

% =====================================================================
\section{Results}
\label{sec:results}
% =====================================================================

\subsection{Decomposing the Protocol Effect}
\label{sec:protocol}
Table~\ref{tab:protocol} reports one system under all six protocols. The
architecture is held fixed at token cross-attention throughout, that being the
configuration in which the ladder was executed; by~\eqref{eq:deltaomega} the
choice of anchor changes the measured effect by at most 2.03 points, and
re-running the ladder under the adopted gated operator is queued
(\TORUN{ladder under gated fusion, for symmetry with Table~\ref{tab:fusion}}).

\begin{table}[!t]
\centering\small
\setlength{\tabcolsep}{3.5pt}
\caption{Effect of evaluation protocol---architecture (token cross-attention),
features, optimiser, schedule and seed are identical across rows; only the
partitioning and the selection partition differ. Intervals for P1 and P3 are
percentile bootstrap 95\% CIs over 2000 resamples of the test partition; P4
reports mean $\pm$ SD across folds. The P1 interval is given for symmetry only:
because that test partition contains exact duplicates of training items, it does
not admit an inferential interpretation.}
\label{tab:protocol}
\begin{tabular}{@{}llccc@{}}
\toprule
& \textbf{Protocol} & \textbf{Acc.\ (\%)} & \textbf{Macro-$F_1$} &
\textbf{UAR} \\
\midrule
P1  & duplicated + random    & 94.79 & 94.97 & 95.13 \\
    & {\footnotesize 95\% CI} & {\footnotesize [92.9, 96.5]} &
      {\footnotesize [93.1, 96.6]} & \\
P2  & deduplicated + random  & \RUN & \RUN & \RUN \\
P2t & \quad + test selection & \RUN & \RUN & \RUN \\
P3  & actor-disjoint, single & 72.50 & 70.83 & 73.05 \\
    & {\footnotesize 95\% CI} & {\footnotesize [67.1, 78.3]} &
      {\footnotesize [64.8, 76.2]} & \\
P3t & \quad + test selection & \RUN & \RUN & \RUN \\
P4  & actor-disjoint, 5-fold & 64.28 & 63.42 & 64.52 \\
    & {\footnotesize $\pm$ SD} & {\footnotesize $\pm$ 5.30} &
      {\footnotesize $\pm$ 4.45} & {\footnotesize $\pm$ 4.59} \\
    & {\footnotesize 95\% CI of mean} & {\footnotesize [57.7, 70.9]} &
      {\footnotesize [57.9, 68.9]} & \\
\midrule
    & \textbf{P1 $-$ P4 (total)} & \textbf{30.51} & \textbf{31.56} &
      \textbf{30.61} \\
    & \textbf{P2 $-$ P4 (attributable)} & \RUN & \RUN & \RUN \\
\bottomrule
\end{tabular}
\end{table}

\begin{figure}[!t]
\centering
\paperfig{figures/fig2_decomposition_ladder}{\figw}{0.735}
\caption{Decomposition of the protocol effect. Each rung of
Table~\ref{tab:protocol} is a step, annotated with the increment attributable to
the single design choice that separates it from the rung below. The three
increments a published system can plausibly incur are solid;
$\Delta_{\mathrm{dup}}$, which arises from a corrupted local copy rather than a
protocol choice, is shown separately. Target file:
\texttt{figures/fig2\_decomposition\_ladder.pdf}.}
\label{fig:ladder}
\end{figure}

\emph{Reading Fig.~\ref{fig:ladder}.} The vertical axis is test accuracy on a
single scale, so the height of each step is the cost of one design decision in
the units the field actually reports. The figure should be read from the top
down: each descent is what a reader loses by removing one source of optimism,
and the total descent is the distance between a number that would pass without
comment in most venues and a number a practitioner could rely on. The separation
of $\Delta_{\mathrm{dup}}$ from the remaining three is deliberate---only the
lower three are choices the published literature makes.

Three observations follow.

First, the contaminated configuration yields 94.79\%. Nothing about this figure
appears anomalous in isolation: it sits among the 95.83\% of Ibrahim \emph{et
al.} \cite{ibrahim2026multimodal} and the 97.8--99.0\% of Jafarzadeh \emph{et
al.} \cite{jafarzadeh2024ssl}, and would be reportable without comment. The same
system under speaker-disjoint five-fold cross-validation attains
64.28\%~$\pm$~5.30---a gap of 30.5 points obtained without changing a line of
the model.

We separate that total into two parts, because they carry different weight as
criticism of the literature. $\Delta_{\mathrm{dup}}$ arises from a defect in our
own extraction pipeline and not from any practice we can attribute to published
work; we report it because it is silent and easy to introduce, not to indict
anyone. The remaining three increments---speaker-dependent partitioning,
test-set model selection, single-split reporting---are choices that appear in
the published record, and their sum $\mu_{\mathrm{P2}}-\mu_{\mathrm{P4}}$ is the
quantity that speaks to the literature.

Second, the single-split speaker-independent figure (72.50\%) exceeds the
cross-validated mean by 8.2 points and lies outside the latter's standard
deviation. This is a sampling effect, not an inconsistency: with four held-out
actors the estimate depends heavily on which four. Per-actor accuracies in that
split range from 58.3\% to 85.0\%, a spread of 26.7 points. We therefore report
P4 as our primary result and regard single-split speaker-independent figures as
uninterpretable without the fold assignment.

Third, the 95\% confidence interval of the P4 mean spans $[57.7, 70.9]$---a
width of 13 points on the quantity the field treats as a single number. That
width is itself an argument.

The magnitude is consistent with the leakage literature outside speech
\cite{rosenblatt2024leakage,elangovan2021memorization}. What is specific to
acted emotional speech is that the leakage is not accidental: it follows
directly from an elicitation design that holds lexical content constant so that
emotion is the only varying factor. The property that makes the corpus a clean
stimulus set is the property that makes random partitioning catastrophic.

\subsection{Fusion Architecture Ablation}
\label{sec:fusion}
Table~\ref{tab:fusion} compares the three operators under P4 on identical folds,
and Table~\ref{tab:perfold} gives the per-fold values from which every
aggregate is computed.

\begin{table}[!t]
\centering\small
\setlength{\tabcolsep}{3pt}
\caption{Fusion-operator ablation under speaker-independent five-fold
cross-validation (P4), identical folds, features, encoder depth,
hyperparameters and head capacity. Mean $\pm$ SD across folds; 95\% CI of the
mean in brackets. No value is bolded: the top two operators are statistically
indistinguishable (Table~\ref{tab:fusion-stats}).}
\label{tab:fusion}
\adjustbox{max width=\linewidth}{%
\begin{tabular}{@{}lcccr@{}}
\toprule
\textbf{Fusion} & \textbf{Acc.\ (\%)} & \textbf{Macro-$F_1$} & \textbf{UAR} &
\textbf{Params} \\
\midrule
Concatenation & 66.32 $\pm$ 5.80 & 65.28 $\pm$ 5.30 & 66.31 $\pm$ 5.74 &
187{,}656 \\
 & {\footnotesize [59.1, 73.5]} & {\footnotesize [58.7, 71.9]} &
   {\footnotesize [59.2, 73.4]} & \\
Gated \emph{(adopted)} & 66.25 $\pm$ 4.97 & 65.38 $\pm$ 4.67 &
65.78 $\pm$ 4.64 & 187{,}656 \\
 & {\footnotesize [60.1, 72.4]} & {\footnotesize [59.6, 71.2]} &
   {\footnotesize [60.0, 71.5]} & \\
Cross-attn.\ ($K{=}4$) & 64.28 $\pm$ 5.30 & 63.42 $\pm$ 4.45 &
64.52 $\pm$ 4.59 & 584{,}456 \\
 & {\footnotesize [57.7, 70.9]} & {\footnotesize [57.9, 68.9]} &
   {\footnotesize [58.8, 70.2]} & \\
\bottomrule
\end{tabular}}
\end{table}

\begin{table}[!t]
\centering\small
\caption{Per-fold test accuracy (\%) under P4. Folds are identical across
operators (Table~\ref{tab:folds}); differences are therefore paired. Fold~4,
whose held-out actors are 2, 10, 12, 13 and 15, is the hardest for all three
operators, which is why the paired design is essential and a single split is
not.}
\label{tab:perfold}
\begin{tabular}{@{}lccccc@{}}
\toprule
\textbf{Fusion} & \textbf{F1} & \textbf{F2} & \textbf{F3} & \textbf{F4} &
\textbf{F5} \\
\midrule
Concatenation          & 73.00 & 69.33 & 63.33 & 58.00 & 67.92 \\
Gated                  & 70.00 & 70.33 & 66.67 & 58.00 & 66.25 \\
Cross-attention        & 67.00 & 69.67 & 63.67 & 55.67 & 65.42 \\
\bottomrule
\end{tabular}
\end{table}

\begin{table}[!t]
\centering\footnotesize
\setlength{\tabcolsep}{2.6pt}
\caption{Pairwise comparisons across the five actor-disjoint folds of P4.
$\Delta$ is the mean paired difference from unrounded fold values;
$d_z = \Delta/\mathrm{SD}(\Delta)$. Holm--Bonferroni is applied within each
metric across its three comparisons. Macro-$F_1$ is the primary metric.}
\label{tab:fusion-stats}
\adjustbox{max width=\linewidth}{%
\begin{tabular}{@{}llccccc@{}}
\toprule
\textbf{Comparison} & \textbf{Metric} & $\bm{\Delta}$ & \textbf{SD} &
\textbf{95\% CI} & $\bm{p_{\mathrm{Holm}}}$ & $\bm{d_z}$ \\
\midrule
Gated $-$ xattn  & Macro-$F_1$ & $+1.97$ & 0.76 & [$+1.02,+2.91$] &
\textbf{0.013} & 2.59 \\
Concat $-$ xattn & Macro-$F_1$ & $+1.86$ & 2.52 & [$-1.27,+4.99$] & 0.348 &
0.74 \\
Gated $-$ concat & Macro-$F_1$ & $+0.10$ & 3.02 & [$-3.64,+3.85$] & 0.942 &
0.03 \\
\midrule
Gated $-$ xattn  & Accuracy & $+1.97$ & 1.15 & [$+0.55,+3.39$] & 0.055 & 1.72 \\
Concat $-$ xattn & Accuracy & $+2.03$ & 2.61 & [$-1.21,+5.27$] & 0.313 & 0.78 \\
Gated $-$ concat & Accuracy & $-0.07$ & 2.44 & [$-3.10,+2.97$] & 0.954 &
$-0.03$ \\
\midrule
Gated $-$ xattn  & UAR & $+1.27$ & 1.02 & [$+0.00,+2.53$] & 0.150 & 1.24 \\
Concat $-$ xattn & UAR & $+1.80$ & 2.59 & [$-1.41,+5.01$] & 0.390 & 0.70 \\
Gated $-$ concat & UAR & $-0.53$ & 3.39 & [$-4.74,+3.67$] & 0.744 & $-0.16$ \\
\bottomrule
\end{tabular}}
\end{table}

Cross-attention is the weakest of the three on all three metrics. Only one
comparison survives multiplicity correction: gated fusion exceeds
cross-attention on the primary metric by 1.97 macro-$F_1$ points
($p_{\mathrm{Holm}} = 0.013$, $d_z = 2.59$). The corresponding accuracy and UAR
comparisons do not survive ($0.055$, $0.150$), and neither concatenation
comparison approaches significance. Gated fusion and concatenation are
indistinguishable, and the sign of their difference flips between metrics
($-0.07$ accuracy, $+0.10$ macro-$F_1$, $-0.53$ UAR)---itself diagnostic that
the comparison is noise.

We adopt gated fusion as the reference configuration, and stress that the choice
is \emph{not empirically forced}: the two are indistinguishable on every metric
($p_{\mathrm{Holm}} \geq 0.74$), and selecting between them on a tenth of a
point would be an instance of the researcher degree of freedom this paper
criticises. We prefer the gate for two architectural reasons rather than a
performance one. It produces a fixed-width representation independent of the
number of streams, so the framework extends without altering downstream
capacity; and by Proposition~\ref{prop:hull} its gate vector is a directly
readable per-coordinate allocation between the two views, which concatenation
does not expose. A reader who prefers concatenation for its lower cost is
equally justified by our results.

\emph{Protocol versus architecture.} The largest architectural difference is
2.03 accuracy points against a 30.5-point protocol effect measured with the
architecture held fixed---a ratio of roughly fifteen to one.
Fig.~\ref{fig:pva} shows both on a common vertical scale.

\begin{figure}[!t]
\centering
\paperfig{figures/fig4_protocol_vs_architecture}{\figw}{0.62}
\caption{Evaluation protocol (left) and fusion architecture (right) on a common
vertical scale. Varying only the partitioning within one fixed system produces a
30.5-point spread; varying the fusion operator at the most defensible protocol
produces 2.0 points. Error bars are 95\% confidence intervals (bootstrap over
test utterances for P1 and P3; of the mean across folds for P4 and the three
operators).}
\label{fig:pva}
\end{figure}

\emph{Reading Fig.~\ref{fig:pva}.} The two panels share axis limits; that shared
scale is the entire content of the figure and nothing should be read from either
panel alone. The left panel spans most of the available range. The right panel
collapses to near-flatness, and the error bars of its three operators overlap
one another almost completely---so the visual impression that the operators are
interchangeable is not an artefact of the scale but the finding itself. A reader
who wishes to see the right panel resolved must either accept a scale on which
the left panel no longer fits, or collect enough folds to shrink the intervals;
Section~\ref{sec:mde} quantifies the second option.

\subsection{Stream Ablation and the Value of the Psychoacoustic View}
\label{sec:modality}
Table~\ref{tab:modality} isolates each stream, and is the direct test of $U_p$.
Unimodal configurations retain identical encoder depth, latent width and head
capacity. We additionally report eGeMAPS \cite{eyben2016gemaps}, extracted with
openSMILE \cite{eyben2010opensmile}, so the perceptual stream is benchmarked
against a standardised parameter set rather than assumed adequate.

\begin{table}[!t]
\centering\small
\setlength{\tabcolsep}{3pt}
\caption{Stream ablation under P4. Identical folds, encoder depth, latent width
$d=128$ and head ($17{,}544$ parameters); only the input stream differs.
Parameter counts are per row because encoder depth does not equalise capacity
across input dimensionalities ($34\!\to\!128$ against $768\!\to\!128$).}
\label{tab:modality}
\adjustbox{max width=\linewidth}{%
\begin{tabular}{@{}lccccr@{}}
\toprule
\textbf{Input} & \textbf{Dim} & \textbf{Acc.} & \textbf{Macro-$F_1$} &
\textbf{UAR} & \textbf{Params} \\
\midrule
Chance                   & ---    & 12.50 & --- & 12.50 & 0 \\
Psychoacoustic only      & 34     & \RUN & \RUN & \RUN & 39{,}304 \\
eGeMAPS only             & 88     & \RUN & \RUN & \RUN & \RUN \\
HuBERT only              & 768    & \RUN & \RUN & \RUN & 133{,}256 \\
\midrule
eGeMAPS + HuBERT (gated) & 88+768 & \RUN & \RUN & \RUN & \RUN \\
Both (gated fusion)      & 34+768 & 66.25 & 65.38 & 65.78 & 187{,}656 \\
\bottomrule
\end{tabular}}
\end{table}

The HuBERT-only row is our in-harness frozen-HuBERT baseline
(Section~\ref{sec:benchmark-method}), so any comparison we make on our own
behalf is against a system we ran on our own folds.

\emph{Complementarity beyond the mean.} Section~\ref{sec:mde} shows this
evaluation cannot resolve differences below 3.47 accuracy points, so a
psychoacoustic contribution smaller than that is invisible in
Table~\ref{tab:modality} whether or not $U_p > 0$.
Algorithm~\ref{alg:complement} therefore examines the error sets directly.
Fig.~\ref{fig:complement} reports the result.
\TORUN{populate Fig.~\ref{fig:complement}; state $J$, $|E_h\setminus E_p|$,
$|E_p\setminus E_h|$ and $p_{\mathrm{McN}}$ here}

\begin{figure*}[!t]
\centering
\paperfig{figures/fig5_stream_complementarity}{\wfigw}{0.350}
\caption{Complementarity of the two streams under P4, on pooled predictions in
which every utterance is classified exactly once by a model that never observed
its speaker. \emph{Left:} overlap of the error sets of the two unimodal systems,
with the counts each stream alone recovers---the operational form of $U_p>0$.
\emph{Centre:} per-class recall of the fused system minus the stronger unimodal
system, so gains offset by losses are visible where aggregate accuracy hides
them. \emph{Right:} distribution of the learned gate $\mathbf{g}$ across the 128
dimensions and the eight categories. Target file:
\texttt{figures/fig5\_stream\_complementarity.pdf}.}
\label{fig:complement}
\end{figure*}

\emph{Reading Fig.~\ref{fig:complement}.} The left panel is the panel that
answers Objective~1, and the quantity to look at is the two crescents rather
than the intersection: a large intersection with empty crescents means the views
fail on the same utterances and the perceptual stream is redundant, while
populated crescents mean $U_p>0$ regardless of what the accuracy column says.
The centre panel guards against a specific misreading of a null aggregate
result---a stream that raises recall on two categories and lowers it on two
others produces zero net accuracy while changing the model substantially. The
right panel is diagnostic rather than evaluative: a gate concentrated near zero
means the model has learned to ignore the perceptual stream, which would corroborate redundancy from the model's own parameters rather than from its outputs.

\subsection{What This Evaluation Can Resolve}
\label{sec:mde}
Applying~\eqref{eq:mde} with the observed standard deviation of paired fold-wise
accuracy differences, $\overline{\mathrm{SD}(\Delta)} = 2.07$ points averaged
over the three accuracy comparisons of Table~\ref{tab:fusion-stats}, a two-sided
paired $t$-test at $\alpha = 0.05$ attains 80\% power only for
$|d_z| \geq 1.68$, giving $\delta_{\min}(5) = \mathbf{3.47}$ accuracy points.
The same calculation gives $2.06$ points at ten folds and $1.23$ points under
leave-one-speaker-out with 24 folds (Fig.~\ref{fig:mde}).

\begin{figure}[!t]
\centering
\paperfig{figures/fig9_minimum_detectable_effect}{0.94\figw}{0.70}
\caption{Design sensitivity as a function of fold count, at 80\% power and
$\alpha = 0.05$, using the observed standard deviation of paired fold-wise
accuracy differences. At five folds the design resolves 3.47 accuracy points.
The largest difference between any two fusion operators (2.03 points, dashed)
and the difference between the top two (0.07 points, dotted) both fall below it.}
\label{fig:mde}
\end{figure}

\emph{Reading Fig.~\ref{fig:mde}.} The curve falls as $1/\sqrt{n}$ modulated by
the $t$ critical value, so the steep early portion is where additional folds buy
the most resolution: moving from five to ten folds recovers 1.4 points of
sensitivity, while moving from ten to twenty-four recovers only 0.8. The two
horizontal reference lines are our own measured effects. The figure should be
read as a design tool rather than a result---given a difference one hopes to
detect, it returns the number of folds required, and given a number of folds it
returns the smallest claim that would be admissible under
Definition~\ref{def:admissible}.

By Definition~\ref{def:admissible}, the difference between the best and worst
fusion operators (2.03 points) and between the top two (0.07 points) are both
\textbf{inadmissible} under our own design. This is the operational content of
the paper's central claim: not merely that architectural effects are smaller
than protocol effects, but that on this corpus, under a protocol a careful
author would defend, they are smaller than the evaluation can measure.

One apparent tension should be addressed, since it is easily misread.
Table~\ref{tab:fusion-stats} reports a 1.97-point macro-$F_1$ difference as
significant after correction, below the 3.47-point threshold. There is no
contradiction: the threshold uses the \emph{average} paired standard deviation
across comparisons (2.07 points), whereas that comparison exhibited an unusually
small one (0.76 points). Design sensitivity characterises what an evaluation can
be relied upon to detect in general, not what it detected in one favourable
case---and with $n=5$ an SD estimate of 0.76 is itself imprecise. We report the
significant comparison and decline to build on it.

\subsection{Per-Class Performance and Confusions}
\label{sec:perclass}
Table~\ref{tab:perclass} contrasts per-class results under the contaminated and
speaker-independent protocols for the same system.

\begin{table}[!t]
\centering\small
\setlength{\tabcolsep}{3.4pt}
\caption{Per-class $F_1$ under the contaminated protocol P1 ($n=576$) and the
speaker-independent single split P3 ($n=240$, actors 2, 11, 16, 23), same
architecture. $\Delta$ is the amount by which contamination flatters each class.
The failure that dominates speaker-independent performance, \emph{happy}, is the
class contamination conceals most.}
\label{tab:perclass}
\begin{tabular}{@{}lcccccc@{}}
\toprule
& \multicolumn{2}{c}{\textbf{P1 (contaminated)}} &
  \multicolumn{3}{c}{\textbf{P3 (speaker-disjoint)}} & \\
\cmidrule(lr){2-3}\cmidrule(lr){4-6}
\textbf{Emotion} & \textbf{Rec.} & $\bm{F_1}$ & \textbf{Prec.} &
\textbf{Rec.} & $\bm{F_1}$ & $\bm{\Delta F_1}$ \\
\midrule
Angry     & 1.000 & 0.950 & 0.692 & 0.844 & 0.761 & $+0.189$ \\
Calm      & 0.961 & 0.980 & 0.846 & 0.688 & 0.759 & $+0.221$ \\
Disgust   & 0.935 & 0.966 & 0.824 & 0.875 & 0.848 & $+0.118$ \\
Fearful   & 0.922 & 0.922 & 0.867 & 0.812 & 0.839 & $+0.083$ \\
Happy     & 0.909 & 0.915 & 0.800 & 0.375 & 0.511 & $\bm{+0.404}$ \\
Neutral   & 1.000 & 0.974 & 0.448 & 0.812 & 0.578 & $+0.396$ \\
Sad       & 0.948 & 0.930 & 0.696 & 0.500 & 0.582 & $+0.348$ \\
Surprised & 0.935 & 0.960 & 0.682 & 0.938 & 0.789 & $+0.171$ \\
\midrule
Macro     & 0.951 & 0.950 & 0.732 & 0.730 & 0.708 & $+0.242$ \\
\bottomrule
\end{tabular}
\end{table}

\emph{Happy} is the dominant failure mode under P3, with recall 0.375; its
errors distribute into \emph{angry} (8 of 32), \emph{surprised} (6) and
\emph{neutral} (5). \emph{Sad} has the second-lowest recall (0.500).
\emph{Neutral} shows high recall but low precision (0.448), indicating that the
class-weighted objective over-predicts the minority class.
Section~\ref{sec:why-fail} accounts for both. Class supports are small and the
intervals correspondingly wide: recall of 0.375 on 32 utterances carries a
binomial 95\% interval of roughly $[0.21, 0.57]$.
\TORUN{recompute the P3 column from pooled P4 predictions, so every utterance is
classified once by a model that never observed its speaker}

The $\Delta F_1$ column carries a finding of its own. Contamination does not
inflate the classes uniformly: it flatters \emph{happy} by 0.404 and
\emph{neutral} by 0.396 while flattering \emph{fearful} by only 0.083. The two
classes a speaker-independent system most needs to fix are precisely the two a
contaminated evaluation reports as solved.

\begin{figure}[!t]
\centering
\paperfig{figures/fig8_confusion_matrices}{\figw}{0.45}
\caption{Row-normalised confusion matrices under (a) the contaminated
random-split protocol P1 and (b) the speaker-independent protocol P3. Classes
are ordered by arousal and valence rather than alphabetically.}
\label{fig:cm}
\end{figure}

\emph{Reading Fig.~\ref{fig:cm}.} The two panels are the same system under two
protocols, so every difference between them is caused by the evaluation design
and none by the model. Panel~(a) is what near-perfect diagonal structure looks
like when a classifier can identify the talker rather than the emotion, and it
is worth dwelling on how unremarkable it appears. Panel~(b) is informative in
its off-diagonal structure: because the classes are ordered by arousal and
valence, the errors form a visible block among the high-arousal states of
opposing valence---\emph{happy}, \emph{angry} and \emph{surprised}---which is
the acoustically predicted confusion and is exactly what
Section~\ref{sec:why-fail} attributes to mean-pooling and to the absence of
voice-quality descriptors. The block is absent from panel~(a) not because the
model has solved it but because it never had to.

\subsection{Contamination Conceals Speaker and Gender Disparity}
\label{sec:speaker}
Under P3, per-actor accuracy ranges from 58.3\% (actor 23) to 85.0\% (actor 2),
and accuracy is 83.3\% on female speakers against 61.7\% on male---a 21.6-point
disparity, comparable to asymmetries reported in fairness analyses of
transformer-based SER \cite{wagner2023dawn}.

The same system evaluated under P1 tells a different story
(Table~\ref{tab:fairness}). The gender gap shrinks to 2.5 points and the
per-actor spread narrows from 26.7 to 14.3 points across all 24 actors. A
fairness audit conducted on the contaminated protocol would conclude the system
is close to parity. It is not.

\begin{table}[!t]
\centering\small
\caption{Contamination conceals the fairness failure. Same architecture, same
features, same seed; only the evaluation protocol differs. A disparity audit run
under P1 would report near-parity for a system that is 21.6 points worse on male
speakers.}
\label{tab:fairness}
\adjustbox{max width=\linewidth}{%
\begin{tabular}{@{}lccc@{}}
\toprule
& \textbf{P1 (contaminated)} & \textbf{P3 (speaker-disjoint)} &
\textbf{Ratio} \\
\midrule
Female accuracy      & 96.0 & 83.3 & --- \\
Male accuracy        & 93.5 & 61.7 & --- \\
\textbf{Gender gap}  & \textbf{2.5} & \textbf{21.6} & $8.6\times$ \\
\midrule
Per-actor min        & 85.7 & 58.3 & --- \\
Per-actor max        & 100.0 & 85.0 & --- \\
\textbf{Actor spread} & \textbf{14.3} & \textbf{26.7} & $1.9\times$ \\
\midrule
Worst-class $F_1$    & 0.915 (\emph{happy}) & 0.511 (\emph{happy}) &
$1.8\times$ \\
\bottomrule
\end{tabular}}
\end{table}

This extends the paper's argument beyond reported accuracy. A protocol that
inflates a headline number by 30 points also suppresses the two audits---per-group disparity and per-class failure---on which responsible deployment depends,
and suppresses them by a larger relative factor than it inflates the headline.
We report the disparity rather than correcting for it; per-speaker normalisation
is an obvious mitigation. Two cautions apply: the P3 figure rests on four
held-out actors and should be recomputed across all 24 from pooled P4
predictions (\TORUN{per-actor and per-gender accuracy, pooled P4}), and the
magnitude of the variation is itself a further argument for cross-validated
reporting.

\subsection{Empirical Confirmation of the Degeneracy Result}
\label{sec:degeneracy}
Proposition~\ref{prop:degen} and Corollary~\ref{cor:inert} predict that the
framework as originally specified---cross-attention applied directly to pooled
stream vectors, so that each sequence has length one---is exactly
gradient-inert. We confirmed this on the original implementation. Scaling the
psychoacoustic input by a factor of 100 while holding the HuBERT input fixed
changed the output logits by $0.0$ to machine precision, and the accumulated
gradient with respect to every parameter of the psychoacoustic encoder was
exactly zero.

Two consequences deserve emphasis. First, any attention map derived from such a
configuration is uninformative: every weight is $1.0$ by construction, and
interpreting it as evidence of learned cross-stream interaction is
unwarranted---a specific instance of the concern raised by Bibal \emph{et al.}
\cite{bibal2022attention}. Second, the failure is silent. The model trains, the
loss decreases, accuracy is competitive, and no diagnostic short of a gradient
check reveals that half the architecture is inert. The assertion at line~11 of
Algorithm~\ref{alg:unified} exists to make it loud.

Correcting the mechanism by expanding each stream into $K = 4$ tokens restores
gradient flow to both encoders. It does not make cross-attention competitive:
Table~\ref{tab:fusion} shows the corrected mechanism underperforming a single
sigmoid gate at $3.1\times$ the parameters and $5.3\times$ the multiply--accumulate cost. We have not swept $K$; the claim should be read as holding at
$K=4$ rather than for all $K$.

\subsection{Computational Complexity}
\label{sec:complexity}
Table~\ref{tab:complexity} instantiates Proposition~\ref{prop:cost}. Counts are
exact, derived analytically and confirmed against the implementation.

\begin{table}[!t]
\centering\small
\setlength{\tabcolsep}{3pt}
\caption{Computational profile of the trainable stack, excluding the frozen
HuBERT encoder. MACs are multiply--accumulate operations per utterance, counting
linear projections and attention products. The trainable stack is negligible
beside one frozen-encoder forward pass ($\approx$16 GMAC for a 3.7~s utterance),
so operator choice is a memory and parameter question rather than a throughput
one.}
\label{tab:complexity}
\begin{tabular}{@{}lrrrr@{}}
\toprule
\textbf{Configuration} & \textbf{Fusion} & \textbf{Total} & \textbf{MACs} &
\textbf{Rel.} \\
 & \textbf{params} & \textbf{params} & & \\
\midrule
Psychoacoustic only    & ---      & 39{,}304  & 38{,}144  & 0.21$\times$ \\
HuBERT only            & ---      & 133{,}256 & 132{,}096 & 0.71$\times$ \\
Concatenation          & 32{,}896 & 187{,}656 & 185{,}600 & 1.00$\times$ \\
Gated \emph{(adopted)} & 32{,}896 & 187{,}656 & 185{,}600 & 1.00$\times$ \\
Cross-attn.\ ($K{=}4$) & 429{,}696 & 584{,}456 & 980{,}224 & 5.28$\times$ \\
\bottomrule
\end{tabular}
\end{table}

Three points follow. The gate costs exactly what concatenation costs, so
Proposition~\ref{prop:hull}'s interpretability advantage is free. Token
cross-attention costs $3.11\times$ the parameters and $5.28\times$ the
operations for the worst accuracy of the three, which is the strongest practical
statement the paper can make about it. And because the frozen encoder dominates
by four to five orders of magnitude, none of these differences is a deployment
consideration---the case for the gate rests on parameter economy and
interpretability, not on latency, and we decline to argue otherwise.

\subsection{Performance Matrix and Benchmarking}
\label{sec:literature}
Table~\ref{tab:matrix} consolidates every system we ran and every published
figure we compare against, under the protocol-stratified discipline of
Section~\ref{sec:benchmark-method}.

\begin{table*}[!t]
\centering\small
\setlength{\tabcolsep}{4pt}
\caption{Performance matrix. Figures are \textbf{not} comparable across protocol
classes, and within a class remain only loosely comparable because encoder
adaptation, fold count and model-selection practice differ. \emph{Disp.}
records whether the source reports any measure of variability. Rows marked
\emph{in-harness} were run by us, on the folds of Table~\ref{tab:folds}, in the
released code. LSO denotes leave-speaker-out; LOSO leave-one-speaker-out.}
\label{tab:matrix}
\begin{tabular}{@{}llccccll@{}}
\toprule
\textbf{System} & \textbf{Protocol} & \textbf{Acc.\ (\%)} &
\textbf{Macro-$F_1$} & \textbf{UAR} & \textbf{Disp.} & \textbf{Params} &
\textbf{Source} \\
\midrule
\multicolumn{8}{l}{\emph{(a) Speaker-dependent --- random utterance partition.
Not comparable with (b).}} \\
wav2vec2 + classifier      & random 80/20  & 99.0  & --- & --- & \checkmark &
--- & \cite{jafarzadeh2024ssl} \\
HuBERT + classifier        & random 80/20  & 97.8  & --- & --- & \checkmark &
--- & \cite{jafarzadeh2024ssl} \\
Hybrid fusion (A+V)        & random split  & 95.83 & --- & --- & $\times$ &
--- & \cite{ibrahim2026multimodal} \\
PAER-LLM, cross-attn.\ (P1) & duplicated + random & 94.79 & 94.97 & 95.13 &
\checkmark & 584{,}456 & \emph{in-harness} \\
MFCC-40 + random forest    & random split  & 92.36 & --- & --- & $\times$ &
--- & \emph{in-harness}\textsuperscript{a} \\
emotion2vec                & random 10-fold & 82.43 & --- & --- & $\times$ &
--- & \cite{ma2024emotion2vec} \\
\midrule
\multicolumn{8}{l}{\emph{(b) Speaker-independent. Our primary result is
PAER-LLM (P4).}} \\
wav2vec2 (weighted layers) & fixed actor split & --- & --- & 84.3 & \checkmark &
--- & \cite{pepino2021wav2vec} \\
CAMuLeNet                  & 10-fold LSO   & 82.3  & --- & --- & $\times$ &
--- & \cite{goel2024camulenet} \\
XLSR-wav2vec2 + MLP        & subj.-wise 5-fold & 81.82 & --- & --- &
\checkmark & --- & \cite{lunajimenez2022multimodal} \\
Whisper-Medium             & 10-fold LSO   & 71.3  & --- & --- & $\times$ &
--- & \cite{goel2024camulenet} \\
PAER-LLM, single split (P3) & actor-disjoint, 1 split & 72.50 & 70.83 & 73.05 &
\checkmark & 584{,}456 & \emph{in-harness} \\
\textbf{PAER-LLM, concat (P4)} & actor-disj.\ 5-fold & \textbf{66.32} &
\textbf{65.28} & \textbf{66.31} & \checkmark & 187{,}656 & \emph{in-harness} \\
\textbf{PAER-LLM, gated (P4)} & actor-disj.\ 5-fold & \textbf{66.25} &
\textbf{65.38} & \textbf{65.78} & \checkmark & 187{,}656 & \emph{in-harness} \\
WavLM                      & 10-fold LSO   & 65.4  & --- & --- & $\times$ &
--- & \cite{goel2024camulenet} \\
\textbf{PAER-LLM, cross-attn.\ (P4)} & actor-disj.\ 5-fold & \textbf{64.28} &
\textbf{63.42} & \textbf{64.52} & \checkmark & 584{,}456 & \emph{in-harness} \\
HuBERT                     & 10-fold LSO   & 58.7  & --- & --- & $\times$ &
--- & \cite{goel2024camulenet} \\
HuBERT only (ours)         & actor-disj.\ 5-fold & \RUN & \RUN & \RUN &
\checkmark & 133{,}256 & \emph{in-harness} \\
Psychoacoustic only (ours) & actor-disj.\ 5-fold & \RUN & \RUN & \RUN &
\checkmark & 39{,}304 & \emph{in-harness} \\
Hybrid fusion (A+V)        & LOSO          & 48.06 & --- & --- & \checkmark &
--- & \cite{ibrahim2026multimodal} \\
wav2vec2.0                 & 10-fold LSO   & 39.6  & --- & --- & $\times$ &
--- & \cite{goel2024camulenet} \\
\midrule
Chance                     & ---           & 12.50 & --- & 12.50 & --- & 0 &
--- \\
\bottomrule
\end{tabular}
\\[2pt]
{\footnotesize \textsuperscript{a}Our own pre-deep-learning baseline. Its
protocol is a random split; like the published rows of class (a) it is not
comparable with class (b), and we include it to show that a 1990s feature set
under a contaminated protocol outscores every speaker-independent system in the
table.}
\end{table*}

Three features of Table~\ref{tab:matrix} carry the benchmarking argument.

First, the highest score in the table belongs to a random-split system, and the
fifth-highest to MFCC features with a random forest---a pre-deep-learning
baseline that outscores every speaker-independent entry, including systems using
encoders that did not exist when it was designed. Under protocol-blind
comparison one would conclude that SER has regressed for a decade. This is the
clearest available demonstration that cross-protocol tabulation is not merely
imprecise but inverts the ordering.

Second, class (b) spans 44.7 points---wider than class (a)'s 16.6---so protocol
annotation alone does not make figures comparable. Fold count, encoder
adaptation and selection practice still differ, and the two entries in class (b)
that exceed our result by the largest margin \cite{goel2024camulenet,
lunajimenez2022multimodal} both adapt the encoder, which we deliberately do not.

Third, the \emph{Disp.} column is mostly empty. A point estimate without
dispersion cannot support a ranking at the one-to-three-point margins the
literature disputes, which is the quantitative form of the recommendation in
Section~\ref{sec:recommendations}.

% =====================================================================
\section{Discussion}
\label{sec:discussion}
% =====================================================================

\subsection{Overall Discussion}
\label{sec:overall}
Read together, the results describe a field measuring below its own noise floor
and not reporting that it is doing so.

The protocol effect is 30.5 accuracy points; the architecture effect at the most
defensible protocol is 2.03; the design's minimum detectable effect is 3.47.
These three numbers are the paper in miniature. The first exceeds the second by
a factor of fifteen, which means that any comparison across papers with
different partitioning schemes is dominated by a nuisance factor. The third
exceeds the second, which means that even a comparison \emph{within} one careful
harness cannot resolve differences of the size the literature routinely reports
as contributions. Our own framework's contribution is subject to the same bound,
and we have applied it to ourselves in Section~\ref{sec:mde} rather than
exempting the proposed method.

Three further findings sharpen this. The degeneracy result
(Proposition~\ref{prop:degen}) shows that a widely used fusion construction can
be not merely weak but formally inert, undetectably so without a gradient check,
which means some published dual-stream results may describe single-stream
models. The convex-hull result (Proposition~\ref{prop:hull}) shows that the two
operators that perform identically are not the same function class, so their
empirical tie is a statement about the data rather than about the operators. And
the fairness result (Section~\ref{sec:speaker}) shows that contamination
suppresses disparity audits by a larger relative factor ($8.6\times$) than it
inflates headline accuracy ($1.5\times$), so a protocol chosen for convenience
propagates into claims about who a system works for.

What the results do \emph{not} show is that psychoacoustic descriptors are
redundant. Section~\ref{sec:whencontrib} distinguishes $U_p = 0$ from
$U_p > 0$-but-unmeasurable, and this evaluation cannot separate them; that is
precisely why Algorithm~\ref{alg:complement} measures error sets rather than
means. The honest position on Objective~1 is that the framework is specified,
its operators are characterised analytically, its cost is quantified, and the
representational question it was built to answer is currently bounded by the
measurement apparatus rather than by the representation.

\subsection{Implications for the Literature}
\label{sec:recommendations}
If a 30-point swing is obtainable by partitioning alone, a reported accuracy
without a stated protocol carries almost no information. We do not suggest
speaker-dependent evaluation is never appropriate---for a system intended to
adapt to a known enrolled user it may be the relevant condition---but the
condition must be stated, and results obtained under it must not be compared
with speaker-independent figures. The confusion is compounded by ambiguous
terminology: ``leave-one-out'' and ``$k$-fold cross-validation'' are used for
both utterance-level and speaker-level partitioning.

We recommend that publications on acted emotional speech corpora report, at
minimum: (i) the number of unique recordings alongside the number of rows used;
(ii) whether partitioning is speaker-disjoint, with the actor assignment
published; (iii) dispersion across folds, not a point estimate; (iv) which
partition was used for model selection; and (v) the design sensitivity of the
evaluation, so readers can judge whether a reported improvement is measurable.
Items (ii) and (iii) were the norm established in 2009
\cite{schuller2009challenge}. Adopting an externally defined split such as
EmoBox \cite{ma2024emobox} satisfies (i)--(ii) without requiring each author to
invent a protocol.

\subsection{Why the Fusion Behaves as It Does}
\label{sec:why-fusion}
Proposition~\ref{prop:hull} establishes that the gate and concatenation span
distinct function classes; their empirical tie therefore requires explanation
from the data. Attention is a mechanism for \emph{selecting among} elements of a
sequence conditioned on a query, and its advantage arises when there is
structure to select over. After mean-pooling, each stream is a single
fixed-length vector with no temporal correspondence to the other, so there is
nothing to align. The $K$-token expansion manufactures a sequence by linear
projection, but the tokens are deterministic functions of the same pooled vector
and carry no independent information; attention over them can only recombine
what a linear layer already provides, at $5.28\times$ the cost. On this reading
the gate and the concatenation are not merely adequate but sufficient: with two
vectors and no correspondence structure, the useful combination operators are
essentially per-coordinate reweightings, and both span that space from opposite
sides of Proposition~\ref{prop:hull}.

This predicts that cross-attention should become advantageous as soon as genuine
sequences are retained---frame-level HuBERT states aligned against frame-level
psychoacoustic contours---consistent with gains reported for temporally-resolved
attentive pooling \cite{leygue2025attentivepooling}. It also identifies where
the unified representation should be built next: not in a better operator over
pooled vectors, but at the frame level, where the two views have something to
align.

\subsection{Why Particular Emotions Fail}
\label{sec:why-fail}
\emph{Happy} fails for over-determined reasons. Acoustically, happiness in acted
speech is a high-arousal state whose principal correlates---elevated mean $F_0$,
increased $F_0$ range, raised intensity, faster rate---are shared with anger and
surprise; the cues that distinguish them are spectral-balance and voice-quality
measures (spectral tilt, harmonics-to-noise ratio, jitter, shimmer) and temporal
contour shape, none of which survive mean-pooling and none of which appear in
our 34-descriptor set. Representationally, mean-pooling discards the rise--fall
contour that distinguishes surprise from sustained joy. At the corpus level,
RAVDESS's \emph{happy} is a portrayed rather than felt state. We predict the
failure would be substantially reduced by retaining temporal structure and
adding voice-quality descriptors---a prediction the eGeMAPS row of
Table~\ref{tab:modality} partially tests, since eGeMAPS includes exactly those
measures.

\emph{Neutral} is over-predicted as a direct consequence of the
inverse-frequency class weighting adopted to compensate for its halved support.
The alternatives---no weighting, resampling, or a margin-based loss---trade
neutral recall against neutral precision. It is also worth noting that
\emph{neutral} is the acoustic origin of the arousal--valence space, so any
weakly realised emotional colouring falls near it; some part of the low
precision is substantive rather than an artefact of weighting.

\subsection{The Language-Intelligence Layer}
\label{sec:llm-discussion}
This layer is the least evidenced part of the work, and we treat it accordingly.
Conditioning a language model on a posterior rather than a label is a defensible
design, but as implemented it is posterior-formatted prompting with an asserted
threshold, and we make no claim for it.

The experiment that would make it load-bearing is the calibration analysis
already specified in Algorithm~\ref{alg:complement}. Computing expected
calibration error and a reliability diagram on pooled cross-validated
predictions would establish whether the posterior carries information about
correctness at all; conditioning accuracy on the top-1 minus top-2 margin would
allow $\tau$ to be set rather than asserted. Computing the same quantities under
P1 and P4 would test a stronger hypothesis: that a contaminated protocol
produces a model that is not merely more accurate but differently
\emph{calibrated}, and therefore one whose downstream interaction policy is
tuned to a confidence it does not possess. Given that contamination already
distorts the fairness audit by $8.6\times$ (Section~\ref{sec:speaker}), we
regard this as likely. Fig.~\ref{fig:reliability} reserves the space.

\begin{figure}[!t]
\centering
\paperfig{figures/fig7_reliability}{\figw}{0.574}
\caption{Posterior calibration under the contaminated and speaker-independent
protocols. \emph{Left:} reliability diagrams for P1 and P4 with expected
calibration error, on pooled predictions. \emph{Right:} accuracy conditioned on
the top-1 minus top-2 posterior margin, from which the hedging threshold $\tau$
of Algorithm~\ref{alg:prompt} is set rather than asserted. Target file:
\texttt{figures/fig7\_reliability.pdf}.}
\label{fig:reliability}
\end{figure}

\emph{Reading Fig.~\ref{fig:reliability}.} In the left panel the diagonal is
perfect calibration; a curve below it marks over-confidence. The comparison to
make is not between either curve and the diagonal but between the two curves:
if P1 and P4 are displaced differently, then a hedging policy tuned on a
contaminated evaluation will hedge at the wrong confidence in deployment,
independently of how much accuracy the contamination added. The right panel is
the design curve for $\tau$: the value at which accuracy rises steeply is the
margin below which the system should decline to assert a state, and reading it
off the plot replaces the asserted 0.15 with a measured quantity.

What this layer does not do is learn multimodal grounding. Approaches that align
speech representations into the language model's embedding space
\cite{kang2025frozen,kim2024paralinguistics} achieve something categorically
different, and we regard that as the appropriate direction for this component.

\subsection{Practical Implications}
For practitioners, published RAVDESS accuracies are not a guide to deployed
performance on unheard speakers, and the gap is not a matter of a few points. A
system selected on a 97\% benchmark figure may deliver in the mid-sixties on a
new talker, and---by Section~\ref{sec:speaker}---may deliver considerably worse
than that for one gender. For affective interfaces in mental-health triage,
driver-state monitoring or call-centre analytics, the procurement question is
not ``what accuracy does it report'' but ``under what partition, over how many
speakers, with what dispersion, and with what per-group breakdown''. For system
builders, the results argue for spending effort on encoder adaptation and
temporal modelling, where the literature shows double-digit effects, rather than
on fusion-operator design over pooled vectors, where we measure effects below
the noise floor and prove one popular operator inert.

\subsection{What Generalises and What Does Not}
\label{sec:generalise}
The numerical results are RAVDESS-specific. The analytical results of
Section~\ref{sec:properties} are corpus-independent: the degeneracy of
single-token attention is a property of the softmax, the convex-hull
confinement a property of the sigmoid, and the cost ordering a property of the
parameterisation. The \emph{magnitude} of the protocol effect is
corpus-dependent and we expect it near an upper bound here, since RAVDESS's
factorial design maximises the near-duplicate structure available to a random
split; on corpora with lexical variety and more speakers such as CREMA-D
\cite{cao2014cremad} the effect should be smaller, and on smaller corpora such
as EmoDB \cite{burkhardt2005emodb} larger \cite{rosenblatt2024leakage}. The
\emph{ordering}---that protocol effects exceed architectural effects---we expect
to be robust, because architectural effects in this literature are consistently
reported at the 1--3 point scale while speaker-dependence effects are
consistently reported at 20--40 points wherever both protocols are run on one
system \cite{ibrahim2026multimodal,portal2025benchmarking}.

\subsection{Future Work}
\label{sec:future}
Multi-corpus replication of the decomposition on CREMA-D \cite{cao2014cremad},
EmoDB \cite{burkhardt2005emodb} and IEMOCAP \cite{busso2008iemocap} using EmoBox
splits \cite{ma2024emobox}; frame-level rather than pooled fusion with genuine
cross-attention, which Section~\ref{sec:why-fusion} identifies as the setting
where the unified representation has room to improve; middle-layer and
weighted-layer representations \cite{pasad2023comparative,chiu2025probing};
standardised psychoacoustic metrics (ISO 532-1 loudness, sharpness, roughness
\cite{iso532}) in place of Bark energies; per-speaker normalisation to address
the gender disparity; deep LLM integration via encoder alignment
\cite{kang2025frozen,kim2024paralinguistics} with human evaluation; and
extension of the decomposition to naturalistic corpora following the Interspeech
2025 challenge protocol \cite{naini2025challenge}.

% =====================================================================
\section{Limitations}
\label{sec:limitations}
% =====================================================================
\begin{enumerate}
  \item \textbf{Single corpus.} All results are on RAVDESS. Acted speech from 24
    actors reading two fixed sentences is not representative of spontaneous
    emotional expression, and cross-corpus generalisation is considerably weaker
    \cite{portal2025benchmarking}.
  \item \textbf{Seed variance unmeasured in the reported tables.} The results
    derive from a single seed. Initialisation variance is plausibly comparable
    to the architectural differences we report \cite{bouthillier2021variance},
    and the multi-seed grid (\TORUN{five seeds}) is required before
    Table~\ref{tab:fusion-stats} can be treated as settled. If seed variance
    swamps the fusion differences, that strengthens rather than weakens the
    central claim.
  \item \textbf{Five folds.} With $n = 5$ the exact Wilcoxon test cannot reach
    conventional significance and $\delta_{\min} = 3.47$ points.
    Leave-one-speaker-out would reduce this to 1.23 and is standard on this
    corpus.
  \item \textbf{The representational claim is bounded by the evaluation.} The
    contribution of the psychoacoustic stream can be established from means only
    if it exceeds $\delta_{\min}$; below that, Algorithm~\ref{alg:complement}
    rather than Table~\ref{tab:modality} carries the argument.
  \item \textbf{Protocol ladder anchored on one operator.} Table~\ref{tab:protocol}
    is executed under token cross-attention. By~\eqref{eq:deltaomega} the anchor
    shifts the result by at most 2.03 points, but symmetry with
    Table~\ref{tab:fusion} would be preferable.
  \item \textbf{Frozen encoder, final-layer pooling.} We use mean-pooled
    HuBERT-base without fine-tuning, discarding temporal structure and forgoing
    the largest available source of improvement \cite{chen2023exploring}, and we
    pool the final hidden state though middle layers carry more paralinguistic
    information \cite{pasad2023comparative,chiu2025probing}. Our absolute
    figures are a floor.
  \item \textbf{No $K$ sweep.} The conclusion that cross-attention is
    uncompetitive is established at $K = 4$ only.
  \item \textbf{Bark band placement.} Our 24 bands are equal-width over
    $[0, z(f_{\mathrm{Nyq}})]$ rather than the canonical Zwicker bands, and the
    22.05~kHz analysis rate discards the upper 1.14~Bark of the audible range.
  \item \textbf{Gender disparity reported but not addressed}, and computed on
    four held-out actors (Section~\ref{sec:speaker}).
  \item \textbf{Latency and memory unprofiled.} Table~\ref{tab:complexity}
    reports parameters and MACs analytically, not wall-clock or peak memory.
  \item \textbf{Language layer unevaluated.} No human study, automatic metric or
    ablation supports Layer~4, and $\tau$ is asserted rather than tuned.
\end{enumerate}

Several of these---single corpus, single seed, five folds, absent dispersion on
some quantities---are precisely the reporting weaknesses this paper criticises.
We state them explicitly rather than omitting them, which is the standard we
advocate.

% =====================================================================
\section{Conclusion}
\label{sec:conclusion}
% =====================================================================
We proposed PAER-LLM, a five-layer framework that unifies perceptually motivated
psychoacoustic descriptors with frozen self-supervised speech representations
through a learned per-dimension gate, and conditions an instruction-tuned
language model on the resulting posterior rather than on a discarded top-1
label. We formalised the unified-representation objective, the
information-theoretic condition under which the perceptual view can contribute,
the evaluation protocol as an operator, and an admissibility criterion that a
representational claim must satisfy.

Evaluating the framework under six protocols differing only in partitioning and
model selection, we measured a 30.5-point spread in reported accuracy with the
architecture held fixed. Under the same conditions the largest difference
between three fusion architectures was 2.03 points---below the 3.47-point
minimum detectable effect of the evaluation, and therefore inadmissible under
our own criterion. Contamination additionally conceals a 21.6-point gender
disparity, reporting it as 2.5, and conceals the dominant per-class failure
entirely.

We proved that single-token cross-modal attention is analytically degenerate: it
silently reduces a dual-stream model to a single-stream one, and no diagnostic
short of a gradient check reveals it. We further showed that gated fusion and
concatenation span distinct function classes despite performing identically, and
that the corrected multi-token attention costs $3.1\times$ the parameters and
$5.3\times$ the operations for the worst accuracy of the three.

Our speaker-independent result, 66.25\%~$\pm$~4.97 accuracy with frozen
HuBERT-base features and gated fusion, is reported with dispersion, confidence
intervals, per-speaker and per-gender breakdowns and published fold assignments.
We release the evaluation harness, including the contaminated-protocol
reproduction.

The field does not lack architectural proposals. It lacks the evaluation
discipline required to tell which of them work.

% =====================================================================
\section*{Reproducibility Statement}
RAVDESS is publicly available under CC BY-NC-SA 4.0
\cite{livingstone2018ravdess} from \url{https://doi.org/10.5281/zenodo.1188976}.
The complete implementation---feature extraction, the three fusion operators,
the actor-disjoint splitting utilities, the contaminated-protocol reproduction,
the evaluation harness and all configuration files---is available at
\TOADD{repository URL} (archived at Zenodo, DOI \TOADD{DOI}). The repository
contains the actor-to-fold assignment of Table~\ref{tab:folds} as a
machine-readable manifest, the per-fold prediction files from which every table
and figure is generated, the scripts that regenerate them, and the
corpus-integrity utility of Algorithm~\ref{alg:integrity}. The environment is
pinned in \texttt{requirements.txt}.

The numbers are cross-consistent in ways checkable without the code: in
Table~\ref{tab:perclass}, the support-weighted mean recall of the P3 column
equals the P3 accuracy of Table~\ref{tab:protocol} (72.50\%), its macro-averaged
recall equals the P3 UAR (73.05\%), and its macro-averaged $F_1$ equals the P3
macro-$F_1$ (70.83\%), each to two decimal places; and the per-fold values of
Table~\ref{tab:perfold} reproduce every mean, standard deviation and paired
difference in Tables~\ref{tab:fusion} and~\ref{tab:fusion-stats}.

\section*{Acknowledgment}
\TOADD{acknowledgements, or delete this section}

\bibliographystyle{IEEEtran}
\bibliography{references}

\end{document}
